\subsection{数量微分方程线性组的基本积分组}

在本小节及下一小节中，我们将考虑如下情形：E 是复数域 C 上的 n 维向量空间（因而在 R 上是 2n 维的），并且对每个 \(t \in J\) ，\(A(t)\) 都是 E 关于 C 上向量空间结构的一个自同态。于是可以把 \(A(t)\) 与它关于 E 的一个基（在域 C 上）的矩阵 \((a_{ij}(t))\) 等同起来，这时的 \(a_{ij}\) 是 J 上有定义且规则的 \(n^{2}\) 个复值函数；若以 \(x_{j}\) ( \(1 \leqslant j \leqslant n\) ) 表示向量 \(x \in E\) 关于所选基的（复）分量，则线性方程

\[
{\frac {d \mathbf {x}}{d t}} = A (t). \mathbf {x} + \mathbf {b} (t)\tag{2}
\]

仍等价于方程组

\[
\frac {d x _ {i}}{d t} = \sum_ {j = 1} ^ {n} a _ {i j} (t) x _ {j} + b _ {i} (t) \quad (1 \leqslant i \leqslant n).\tag{3}
\]

于是，定理 1 (IV, p.~179) 与定理 2 (IV, p.~179) 以及命题 2 (IV, p.~181) 表明，对 E 内的每个 \(\mathbf{x}_0 = (x_{k0})_{1\leqslant k\leqslant n}\) ，存在一个且仅一个积分 \(\mathbf{u} = (u_k)_{1\leqslant k\leqslant n}\) 满足方程

\[
{\frac {d \mathbf {x}}{d t}} = A (t). \mathbf {x}\tag{4}
\]

它在 E 上定义，并且取值 \(x_{0}\) 于点 \(t_{0}\) ；这个积分可以写成

\[
\mathbf {u} (t, t _ {0}, \mathbf {x} _ {0}) = C (t, t _ {0}). \mathbf {x} _ {0},
\]

其中 \(C(t, t_{0})\) 是一个 n 阶可逆方阵 \((c_{ij}(t, t_{0}))\) ，其元素是 J × J 上的连续复值函数，并且使得 \(t \mapsto c_{ij}(t, t_{0})\) 是 J 上某个规则函数的原函数。

在 \(n = 1\) 的特殊情形，方程组 (3) 化为单个数量方程

\[
{\frac {d x}{d t}} = a (t) x + b (t)\tag{11}
\]

\((a(t) \text{ 与 } b(t) \text{ 为 J 上的复值规则函数}); \text{ 立即可验证，(单元素)矩阵 } C(t, t_0) \text{ 等于 } \exp\left(\int_{t_0}^{t} a(s) ds\right); \text{ (11) 的积分，其值等于 } x_0 \text{，在点 } t_0 \text{ 处，因而明确地由下式给出}\)

\[
u (t) = x _ {0} \exp \left(\int_ {t _ {0}} ^ {t} a (s) d s\right) + \int_ {t _ {0}} ^ {t} b (s) \exp \left(\int_ {t _ {0}} ^ {t} a (\tau) d \tau\right) d s.\tag{12}
\]

在由 J 到 E 的连续映射组成、赋予紧收敛拓扑的空间 \(\mathcal{C}(\mathrm{J};\mathrm{E})\) 中，方程 (4) 的积分所成的集 I 是一个向量子空间（在 C 上），它同构于 E，从而同构于 \(C^{n}\) (IV, p.~180, cor. 1 及 IV, p.~181, prop. 2)。

这个空间（在域 C 上）的一个基 \((\mathbf{u}_{j})_{1\leqslant j\leqslant n}\) 称为 (4) 的一个基本积分组。

命题 4. 为使方程 (4) 的 n 个积分 \(\mathbf{u}_{j}\ (1 \leqslant j \leqslant n)\) 构成一个基本积分组，必须且只需它们的值 \(\mathbf{u}_{j}(t_{0})\) （在某个点 \(t_{0} \in J\) 处取值）是 E 中线性无关的向量。

事实上，把每个 \(\mathbf{x}_0 \in \mathrm{E}\) 映为积分 \(t \mapsto C(t, t_0). \mathbf{x}_0\) 的映射是 \(\mathrm{E}\) 到 \(\mathcal{I}\) 上的一个同构 (IV, p.~180, cor. 1 及 IV, p.~181, prop. 2)。

若 \(\left(\mathbf{e}_j\right)_{1\leqslant j\leqslant n}\) 是 E 在 C 上的任意一个基，则 \(n\) 个积分

\[
\mathbf {u} _ {j} (t) = C \left(t, t _ {0}\right). \mathbf {e} _ {j} \quad (1 \leqslant j \leqslant n)
\]

便构成一个基本积分组；若把 \(C(t, t_0)\) 与它关于基 \((\mathbf{e}_j)\) 的矩阵等同起来，则诸积分 \(\mathbf{u}_j\) 正是矩阵 \(C(t, t_0)\) 的各列。于是，(4) 的取值 \(\mathbf{x}_0 = \sum_{j=1}^{n} \lambda_j \mathbf{e}_j\) 于点 \(t_0\) 的积分就是 \(C(t, t_0). \mathbf{x}_0 = \sum_{k=1}^{n} \lambda_k \mathbf{u}_k(t)\) 。

对 (4) 的任意 \(n\) 个积分 \(\mathbf{u}_j\) ( \(1 \leqslant j \leqslant n\) )，我们把这 \(n\) 个积分在点 \(t \in \mathbf{J}\) 处关于基 \((\mathbf{e}_j)_{1 \leqslant j \leqslant n}\) （在 \(E\) 中）的行列式称为

\[
\Delta (t) = \left(\mathbf {u} _ {1} (t), \mathbf {u} _ {2} (t), \dots , \mathbf {u} _ {n} (t)\right)\tag{13}
\]

即 \(n\) 个向量 \(\mathbf{u}_j(t)\) 关于基 \((\mathbf{e}_j)\) 的行列式 (Alg., III, p.~522)。我们有 (Alg., III, p.~523, prop. 2)

\[
\Delta (t) = \Delta (t _ {0}) \det \left(C (t, t _ {0})\right).\tag{14}
\]

由 IV, p.~184 的命题 4，为使 \(\left(\mathbf{u}_{j}\right)_{1\leqslant j\leqslant n}\) 是 (4) 的一个基本积分组，必须且只需行列式 \(\Delta(t)\) （诸 \(u_{j}\) 的行列式）\(\neq0\) 在 J 的某个点 \(t_{0}\) 处成立；于是公式 (14) 表明，在 J 的每一点处 \(\Delta(t)\neq0\)，换言之，向量 \(\mathbf{u}_{j}(t)\) ( \(1\leqslant j\leqslant n\) ) 恒线性无关。

命题 5. 矩阵 \(C(t, t_{0})\) 的行列式由下列公式给出

\[
\det \left(C (t, t _ {0})\right) = \exp \left(\int_ {t _ {0}} ^ {t} \operatorname{Tr} (A (s)) d s\right).\tag{15}
\]

事实上，若令 \(\delta(t) = \det\left(C(t, t_0)\right)\) ，则由行列式的求导公式 (I, p.~8, 公式 (3)) 有

\[
\frac {d \delta}{d t} = \operatorname{Tr} \left(\frac {d C (t , t _ {0})}{d t} \left(C (t, t _ {0})\right) ^ {- 1}\right) \delta (t)
\]

亦即，由 \(C(t, t_0)\) 所满足的 IV, p.~179 的微分方程 (5)，有

\[
\frac {d \delta}{d t} = \mathrm{Tr} \big (A (t) \big) \delta (t).
\]

由于 \(\delta(t_0) = 1\) ，由数量线性方程的积分的表达式 (12) (IV, p.~183) 即得公式 (15)。

正如刚才所见，给定 (4) 的 n 个线性无关的积分，便确定了这个方程的全部积分。现在我们要证明，当 \(1 \leqslant p \leqslant n\) 时，知道了方程 (4) 的 p 个线性无关的积分 \(\mathbf{u}_{j}\) ( \(1 \leqslant j \leqslant p\) )，便可把这个方程的积分化为由 n - p 个数量方程组成的齐次线性方程组的积分。设在一个区间 \(K \subset J\) 上存在 n - p 个映射

\[
\mathbf {u} _ {p + k} \quad (1 \leqslant k \leqslant n - p)
\]

它们把 K 映入 E，都是 K 上规则函数的原函数，并且使得对每个 \(t \in K\) ，n 个向量 \(\mathbf{u}_{j}(t) (1 \leqslant j \leqslant n)\) 构成 E 的一个基。

对每个点 \(t_1 \in \mathbf{J}\) ，总存在一个区间 \(\mathbf{K}\) ，即 \(t_1\) 在 \(\mathbf{J}\) 中的一个邻域，在其上可以定义具有前述性质的 \(n - p\) 个函数 \(\mathbf{u}_{p+k}\) ( \(1 \leqslant k \leqslant n - p\) )。因为，设 \((\mathbf{e}_i)_{1 \leqslant i \leqslant n}\) 是 \(\mathbf{E}\) 的一个基；该基中存在 \(n - p\) 个向量，它们与诸 \(\mathbf{u}_j(t_1)\) ( \(1 \leqslant j \leqslant p\) ) 一起构成 \(\mathbf{E}\) 的一个基 (Alg., II, p.~292, th. 2)；例如设它们是 \(\mathbf{e}_{p+1}, \ldots, \mathbf{e}_n\) ；由于行列式 \(\det(\mathbf{u}_1(t), \ldots, \mathbf{u}_p(t), \mathbf{e}_{p+1}, \ldots, \mathbf{e}_n)\) （关于基 \((\mathbf{e}_i)\) ）是 \(t\) 的连续函数并且在 \(t = t_1\) 处不为零，故存在一个邻域 \(\mathbf{K}\) 包含 \(t_1\) ，使行列式在其上不为零；于是可取 \(\mathbf{u}_{p+k}(t) = \mathbf{e}_{p+k}\) ( \(1 \leqslant k \leqslant n - p\) ) （对 \(t \in \mathbf{K}\) ）。

存在一个 n 阶可逆矩阵 \(B(t)\) ，其元素是 K 上规则函数的原函数，使得 \(B(t).e_{j} = \mathbf{u}_{j}(t)\) （对 \(1 \leqslant j \leqslant n\) ）。令 \(\mathbf{x} = B(t).y\) ；则 y 满足方程 \(\frac{dB}{dt}.y + B(t).\frac{dy}{dt} = A(t)B(t).y\) ，它也可以写成

\[
\frac {d \mathbf {y}}{d t} = (B (t)) ^ {- 1} \left(A (t) B (t) - \frac {d B}{d t}\right). \mathbf {y} = H (t). \mathbf {y}
\]

其中 \(H(t) = \left(h_{jk}(t)\right)\) 是一个元素在 K 上为规则函数的矩阵。由 \(B(t)\) 的定义，这个线性方程以 p 个常向量 \(\mathbf{e}_{j}\) ( \(1 \leqslant j \leqslant p\) ) 为积分；由此立即推出，必有 \(h_{jk}(t) = 0\) （对 \(1 \leqslant k \leqslant p\) ）；于是 y 的分量 \(y_{k}\) （关于基 \((\mathbf{e}_{i})\) ）中下标 \(k \geqslant p + 1\) 的那些，满足一个由 n - p 个方程组成的齐次线性方程组；一旦确定了这个方程组的解，\(dy_{j}/dt\) （对下标 \(j \leqslant p\) ）便是 \(y_{k}\) （其中 \(k \geqslant p + 1\) ）的线性函数，从而是已知的，而这些函数的原函数便给出 \(y_{j}\) （对下标 \(j \leqslant p\) ）。

特别地，当已知 IV, p.~183 的方程 (4) 的 \(n - 1\) 个线性无关的积分时，这个方程的积分便化为单个齐次数量方程的积分，随后只需计算 \(n\) 个原函数。

注. 1) 上述一切也适用于如下情形：E 是域 R 上的 n 维空间，并且 \(A(t)\) 对每个 \(t \in J\) 都是 E 的一个自同态：只需处处把 C 换成 R。

\begin{enumerate}
\def\labelenumi{\arabic{enumi})}
\setcounter{enumi}{1}
\tightlist
\item
  设 \(A(t) = (a_{ij}(t))\) 是一个 n 阶方阵，其元素是 t 在 J 上的实值（相应地，复值）规则函数，又设 \(C(t, t_0) = (c_{ij}(t, t_0))\) 是相应的线性方程组 (3) (IV, p.~22) 的预解矩阵。设 F 是 R（相应地，C）上的任意一个完备赋范空间，并考虑线性微分方程组
\end{enumerate}

\[
\frac {d \mathbf {y} _ {i}}{d t} = \sum_ {j = 1} ^ {n} a _ {i j} (t) \mathbf {y} _ {j}
\]

其中未知函数 \(\mathbf{y}_j\) 取值于 F。立即可见，该方程组的解 \((\mathbf{u}_j)_{1\leqslant j\leqslant n}\) ——它满足 \(\mathbf{u}_j(t_0) = \mathbf{d}_j\) （对 \(1\leqslant j\leqslant n\) ； \(\mathbf{d}_j\) 在 F 中任意取值）——由下列公式给出

\[
\mathbf {u} _ {i} (t) = \sum_ {j = 1} ^ {n} c _ {i j} (t, t _ {0}) \mathbf {d} _ {j} \quad (1 \leqslant i \leqslant n).
\]

我们特别考虑如下情形：\(A(t)\) 是 C 上 n 维向量空间 E 的一个自同态，并且存在 E 的一个基，使 \(A(t)\) 关于该基的矩阵对每个 \(t \in J\) 都有实元素。于是上面的结果表明（由 IV, p.~179 的定理 1），预解矩阵 \(C(t, t_{0})\) 关于同一个基也有实元素：只需考虑由所考虑的 E 的基在 R 上生成的向量空间 \(E_{0}\) ，并注意到 \(A(t)\) 在 \(E_{0}\) 上的限制是这个向量空间的一个自同态。

\subsection{伴随方程}

仍然假设空间 E 是 C 上的 n 维空间，设 \(E^{*}\) 是它的对偶 (A, II, p.~40)，这是一个 C 上的 n 维空间 (Alg., II, p.~299, th. 4)；典范双线性型 \(\langle x, x^{*}\rangle\) （定义在 \(E \times E^{*}\) 上）(Alg., II, p.~234) 在这个乘积空间上连续（因为它是 \(x \in E\) 与 \(x^{*} \in E^{*}\) 的分量的多项式）。

给定齐次线性方程 (4) (IV, p.~183)，其中 \(t \mapsto A(t)\) 是 \(\mathbf{J}\) 到 \(\mathcal{L}(\mathrm{E})\) 中的一个规则映射，我们要考察是否存在一个映射 \(t \mapsto \mathbf{v}(t)\) ，它把 \(\mathbf{J}\) 映入 \(\mathrm{E}^*\) ，是 \(\mathbf{J}\) 上某个规则函数的原函数，并且使得数量函数 \(t \mapsto \langle \mathbf{u}(t), \mathbf{v}(t) \rangle\) 在 \(\mathbf{J}\) 上为常数（当 \(\mathbf{u}\) 是 (4) 的任意解时）；这等同于写出：这个函数的导数在 \(\mathbf{u}\) 与 \(\mathbf{v}\) 都可导的每一点处为零，也就是说，必须有

\[
\left\langle \frac {d {\bf u}}{d t}, {\bf v} (t) \right\rangle + \left\langle {\bf u} (t), \frac {d {\bf v}}{d t} \right\rangle = 0
\]

在这样的点处。

现在，由 (4) 有 \(\left\langle\frac{d\mathbf{u}}{dt},\mathbf{v}(t)\right\rangle=\langle A(t).\mathbf{u}(t),\mathbf{v}(t)\rangle=-\langle\mathbf{u}(t),B(t).\mathbf{v}(t)\rangle\) ，其中 \(-B(t)\) 是 \(A(t)\) 的转置 (Alg., II, p.~234)。于是 v 所必须满足的关系可以写成

\[
\left\langle \mathbf {u} (t), \frac {d \mathbf {v}}{d t} - B (t). \mathbf {v} (t) \right\rangle = 0
\]

在 \(A(t)\) 连续且 \(\mathbf{v}(t)\) 可导的所有点处。现在，对这样的点 \(t\) 以及任意的点 \(\mathbf{x}_0 \in E\) ，由 IV, p.~179 的定理 1，存在 (4) 的一个解 \(\mathbf{u}\) 使得 \(\mathbf{u}(t) = \mathbf{x}_0\) ；因此必须有 \(\left\langle \mathbf{x}_0, \frac{d\mathbf{v}}{dt} - B(t).\mathbf{v}(t) \right\rangle = 0\) （对所有的 \(\mathbf{x}_0 \in E\) ），这就蕴含 \(\frac{d\mathbf{v}}{dt} - B(t).\mathbf{v}(t) = 0\) 。因此：

命题 6. 映射 \(t \mapsto \mathbf{v}(t)\) （由 \(\mathbf{J}\) 到 \(\mathrm{E}^*\) ，且为 \(\mathbf{J}\) 上某个规则函数的原函数）使得 \(\langle \mathbf{u}(t), \mathbf{v}(t) \rangle\) 在 \(\mathbf{J}\) 上对 IV, p.~183 的方程 (4) 的每个解 \(\mathbf{u}\) 均为常数，当且仅当 \(\mathbf{v}\) 是齐次线性方程

\[
{\frac {d \mathbf {x}}{d t}} = B (t). \mathbf {x}\tag{16}
\]

的解，其中 \(-B(t)\) 是 \(A(t)\) 的转置。

方程 (16) 称为 (4) 的伴随方程；显然 (4) 也是 (16) 的伴随方程。由于矩阵 \(B(t)\) 的元素是 t 在 J 上的规则函数，前面关于线性方程得到的结果都适用于 (16)。特别地，(16) 的取值 \(x_{0}^{*}\) 于点 \(t_{0}\) 的积分可以写成 \(H(t, t_{0})\) . \(x_{0}^{*}\) ，其中 \(H(t, t_{0})\) 是 \(E^{*}\) 到其自身上的一个双射线性映射，它等同于方程

\[
\frac {d V}{d t} = B (t) V\tag{17}
\]

的取值 \(I\) 于点 \(t_0\) 的积分。由此得到（采用 IV, p.~179 的记号）

\[
\left\langle C (t, t _ {0}). \mathbf {x} _ {0}, H (t, t _ {0}). \mathbf {x} _ {0} ^ {*} \right\rangle = \left\langle \mathbf {x} _ {0}, \mathbf {x} _ {0} ^ {*} \right\rangle
\]

（对任意的 \(x_{0} \in E\) 与 \(x_{0}^{*} \in E^{*}\) ），这表明

\[
H (t, t _ {0}) = \check {C} (t, t _ {0})\tag{18}
\]

（即 \(C(t, t_{0})\) 的逆步映射）。特别地，若已知伴随方程 (16) 的一个基本积分组，则矩阵 \(H(t, t_{0})\) 就被确定，从而 \(C(t, t_{0})\) 也随之确定，于是方程 (4) 的全部积分也都随之确定。

注. 设 E 与 F 是 R（或 C）上的两个完备赋范空间，\((\mathbf{x}, \mathbf{y}) \mapsto \langle \mathbf{x}, \mathbf{y} \rangle\) 是 \(E \times F\) 上的一个连续双线性型，使得关系 `` \(\langle x, y \rangle = 0\) 对所有 \(y \in F\) ''（相应地，`` \(\langle x, y \rangle = 0\) 对所有 \(x \in E\) ''）蕴含 x = 0（相应地，y = 0）。再假设对每个 \(t \in J\) 存在一个 F 到其自身的连续线性映射 \(B(t)\) ，使得 \(\langle A(t).x, y \rangle + \langle x, B(t).y \rangle = 0\) 对每个 \((\mathbf{x}, \mathbf{y}) \in E \times F\) 成立。在这些条件下，与前面一样可以看出，为使一个把 J 映入 F 的映射 \(t \mapsto v(t)\) （某个规则函数的原函数）满足 \(\langle u(t), v(t) \rangle\) 对 (4) 的每个积分 u 都为常数，必须且只需 v 是 (16) 的一个积分，我们仍称 (16) 为 (4) 的伴随方程。

\subsection{常系数线性微分方程}

再设 E 是 R 上一个任意的完备赋范空间；设 A 是 A 的一个连续自同态，与 t 无关，并考虑齐次线性方程

\[
{\frac {d \mathbf {x}}{d t}} = A. \mathbf {x}.\tag{19}
\]

当 \(\mathrm{E}\) 为有限维时，方程 (19) 等价于一个由数量微分方程组成的齐次方程组 (3) (IV, p.~183)，其中系数 \(a_{ij}\) 为常数。

由定理 1 (IV, p.~179)，(19) 的每个积分都在整个 \(\mathbf{R}\) 上有定义；由定理 2 (IV, p.~179)，(19) 的取值 \(\mathbf{x}_0\) （在点 \(t_0 \in \mathbf{R}\) 处）的积分可以写成 \(C(t, t_0)\mathbf{x}_0\) ，其中 \(C(t, t_0)\) 是 E 到其自身上的一个双射双连续线性映射，它满足方程

\[
\frac {d U}{d t} = A U\tag{20}
\]

并且使得 \(C(t_0, t_0) = I\) 。此外，我们有恒等式

\[
C (t + \tau , t _ {0} + \tau) = C (t, t _ {0})\tag{21}
\]

（对任意 \(\tau \in \mathbf{R}\) ）：事实上，由 (20) 有 \(dC(s, t_0 + \tau) / ds = AC(s, t_0 + \tau)\) ，并且由于 \(A\) 是常数，还有

\[
\frac {d C (t + \tau , t _ {0} + \tau)}{d t} = A C (t + \tau , t _ {0} + \tau);
\]

此外

\[
C (t _ {0} + \tau , t _ {0} + \tau) = I = C (t _ {0}, t _ {0}),
\]

由此即得恒等式 (21)，因为在点 \(t_{0}\) 处等于 I 的 (20) 的积分是唯一的。

若令 \(C_0(t) = C(t, 0)\) ，则 \(C(t, t_0) = C_0(t - t_0)\) ；此外，对每个 \(\lambda \in \mathbf{R}\) ，\(C_0(\lambda t)\) 都等同于方程

\[
\frac {d U}{d t} = \lambda A U\tag{22}
\]

的取值 \(I\) 于点 0 的积分。我们给出下面的定义：

定义 1. 给定 E 的一个连续自同态 A，我们以 \(e^{A}\) 或 exp A 表示 E 的这样一个自同构：它等于方程 (20) 的在点 t = 0 处取值 I 的积分在点 t = 1 处的值。

采用此记号，定义 1 之前的注记表明

\[
C (t, t _ {0}) = \exp \left(A (t - t _ {0})\right).\tag{23}
\]

刚才引入的指数记号之所以合理，是基于下面这些性质，它们与 z 为实数或复数时的函数 \(\exp z\) 的性质完全类似 (cf.~III, p.~98 与 106)：

命题 7. \(1^{\circ}\) 映射 \(X \mapsto e^{X}\) 是 \(\mathcal{L}(\mathrm{E})\) 到 \(\mathrm{E}\) 的自同构群（\(\mathcal{L}(\mathrm{E})\) 中的可逆元所成的群）的一个连续映射。

\(2^{\circ}\) 映射 \(t\mapsto e^{Xt}\) （由 \(\mathbf{R}\) 到 \(\mathcal{L}(\mathrm{E})\) ）是可导的，并且

\[
\frac {d}{d t} \big (e ^ {X t} \big) = X e ^ {X t} = e ^ {X t} X.\tag{24}
\]

\(3^{\circ}\) 对任意 \(X\in \mathcal{L}(\mathrm{E})\) 有

\[
e ^ {X} = \sum_ {n = 0} ^ {\infty} \frac {X ^ {n}}{n !}\tag{25}
\]

右端在 \(\mathcal{L}(\mathrm{E})\) 的每个有界子集上绝对且一致收敛；特别地，\(e^{It} = e^{t}I\) （对 \(t \in R\) ）。

\(4^{\circ}\) 若 X 与 Y 可交换，则 Y 与 \(e^{Y}\) 都与 \(e^{X}\) 可交换，并且

\[
e ^ {X + Y} = e ^ {X} e ^ {Y}.\tag{26}
\]

关系式 (24) 由 \(C(t,0)\) 的表达式 (23) 以及该函数是 (20) 的积分这一事实得出；对 n 用递归法由 (24) 推出，\(t \mapsto e^{Xt}\) 在 R 上无穷次可导，并且

\[
\mathrm{D} ^ {n} \left(e ^ {X t}\right) = X ^ {n} e ^ {X t}.
\]

于是由泰勒公式可以写出

\[
e ^ {X} = I + \frac {X}{1 !} + \frac {X ^ {2}}{2 !} + \dots + \frac {X ^ {n}}{n !} + X ^ {n + 1} \int_ {0} ^ {1} \frac {(1 - t) ^ {n}}{n !} e ^ {X t} d t.\tag{27}
\]

另一方面，IV, p.~181 的 cor. 3 表明 \(\left\| e^{Xt}\right\| \leqslant \exp \left(\| X\| |t|\right)\) 。于是公式 (27) 中的余项 \(r_n(X) = X^{n + 1}\int_0^1\frac{(1 - t)^n}{n!} e^{Xt}dt\) 满足不等式

\[
\| r _ {n} (X) \| \leqslant \frac {\| X \| ^ {n + 1}}{(n + 1) !} e ^ {\| X \|}
\]

由此推出公式 (25)，右端的级数在 \(\mathcal{L}(\mathrm{E})\) 的每个有界子集上绝对且一致收敛。对 \(\mathcal{L}(\mathrm{E})\) 中每一对元素 X、T 有

\[
e ^ {X + T} - e ^ {X} = \sum_ {n = 1} ^ {\infty} \frac {1}{n !} \big ((X + T) ^ {n} - X ^ {n} \big).
\]

现在，可以写出 \((X + T)^n - X^n = \sum_{(V_i)} V_1 V_2 \ldots V_n\) ，其中求和遍及 \(2^n - 1\) 个序列 \((V_i)\) （由 \(\mathcal{L}(\mathrm{E})\) 的元素组成），它们满足 \(V_i = X\) 或 \(V_i = T\) （对 \(1 \leqslant i \leqslant n\) ），且至少有一个 \(V_i\) 等于 \(T\) ；于是不等式

\[
\left\| (X + T) ^ {n} - X ^ {n} \right\| \leqslant \left(\| X \| + \| T \|\right) ^ {n} - \| X \| ^ {n}
\]

立即成立，由此

\[
\| \exp (X + T) - \exp X \| \leqslant \exp \left(\| X \| + \| T \|\right) - \exp \| X \|\tag{28}
\]

这就建立了映射 \(X \mapsto \exp X\) 的连续性。

最后，若 X 与 Y 可交换，则 Y 与 \(e^{Xt}\) 可交换 (IV, p.~181, prop. 2)，于是

\[
\frac {d}{d t} \big (e ^ {X t} e ^ {Y t} \big) = X e ^ {X t} e ^ {Y t} + e ^ {X t} (Y e ^ {Y t}) = (X + Y) e ^ {X t} e ^ {Y t}.
\]

另一方面，由于 \(e^{Xt}e^{Yt}\) 在 t = 0 处等于 I，故有 \(e^{Xt}e^{Yt} = e^{(X+Y)t}\) ，由此得公式 (26)。由后者特别可推出，对任意的实数 s 与 t 有

\[
e ^ {X (s + t)} = e ^ {X s} e ^ {X t}\tag{29}
\]

并且还有

\[
e ^ {- X} = \left(e ^ {X}\right) ^ {- 1}.\tag{30}
\]

反之，要注意的是：若不再假设 X 与 Y 可交换，(26) 未必仍成立；因为否则它将蕴含 \(\exp X\) 与 \(\exp Y\) 总可交换，而事实并非如此，简单的例子即可表明这一点 (IV, p.~204, exerc. 3)。

现在假设 E 是域 C 上的有限维向量空间，A 是 E 的一个自同态（关于 E 在 C 上的向量空间结构），可把它与 A 关于 E 的某个基的矩阵等同起来；于是对所有的 \(t \in R\) ，\(e^{At}\) 都是 E 关于同一结构的一个自同构 (IV, p.~181, prop. 2)。设 \(r_k\) ( \(1 \leqslant k \leqslant q\) ) 是自同态 A 的特征多项式 \(\varphi(r) = \det(A - rI)\) 的互异根（在 C 中）（即 A 的“特征根”）；若 \(n_k\) 是 \(r_k\) 的重数，则 \(\sum_{k=1}^{q} n_k = n\) 。我们知道 (Alg., VII, 31, n° 3)，对每个根 \(r_k\) 对应着 E 的一个子空间 \(E_k\) ，其维数为 \(n_k\) ，使得 \(E_k\) 在 A 之下稳定，并且 E 是诸 \(E_k\) 的直和：\(E_k\) 可定义为满足下式的向量 x 所成的子空间

\[
(A - r _ {k} I) ^ {n _ {k}}. \mathbf {x} = 0.
\]

设 \(\mathbf{a}\) 是 E 中任意一个向量；可以写出 \(\mathbf{a} = \sum_{k=1}^{q} \mathbf{a}_k\) ，其中 \(\mathbf{a}_k \in \mathrm{E}_k\) ；于是 IV, p.~188 的方程 (19) 的取值 \(\mathbf{a}\) 于点 \(t = 0\) 的积分由下式给出

\[
\mathbf {u} (t) = e ^ {A t}. \mathbf {a} = \sum_ {k = 1} ^ {q} e ^ {A t}. \mathbf {a} _ {k} = \sum_ {k = 1} ^ {q} e ^ {r _ {k} t} e ^ {(A - r _ {k} I) t}. \mathbf {a} _ {k}.
\]

但由于 \(\mathbf{a}_k\in \mathrm{E}_k\) ，有

\[
\begin{array}{c} e ^ {(A - r _ {k} I) t}. \mathbf {a} _ {k} = \mathbf {a} _ {k} + \frac {t}{1 !} (A - r _ {k} I). \mathbf {a} _ {k} + \frac {t ^ {2}}{2 !} (A - r _ {k} I) ^ {2}. \mathbf {a} _ {k} + \dots \\ + \frac {t ^ {n _ {k} - 1}}{(n _ {k} - 1) !} (A - r _ {k} I) ^ {n _ {k} - 1}. \mathbf {a} _ {k}. \end{array}\tag{31}
\]

于是，IV, p.~188 的方程 (19) 的每个积分都可以写成

\[
\mathbf {u} (t) = \sum_ {k = 1} ^ {q} e ^ {r _ {k} t} \mathbf {p} _ {k} (t)\tag{32}
\]

其中 \(\mathbf{p}_{k}(t)\) 是 t 的一个多项式，其系数在向量空间 \(E_{k}\) 中，次数 \(\leqslant n_{k}-1\) 。特别地，若 A 的特征方程的所有根都是单根，则诸空间 \(E_{k}\) ( \(1 \leqslant k \leqslant n\) ) 都是域 C 上的一维空间，于是存在 n 个向量 \(c_{k}\) ，使得 n 个函数 \(e^{r_{k}t}c_{k}\) ( \(1 \leqslant k \leqslant n\) ) 构成 IV, p.~188 的方程 (19) 的一个基本积分组。

自同态 \(A\) 的特征根也称为 IV, p.~188 的线性方程 (19) 的特征根。可以注意到，\(A\) 的特征方程可以这样得到：写出函数 \(\mathbf{c}e^{rt}\) （对一个向量 \(\mathbf{c} \neq 0\) ）是 (19) 的积分。

当根 \(r_{k}\) ( \(1 \leqslant k \leqslant q\) ) 连同 \(n_{k}\) （\(r_{k}\) 的重数）都已明确求出之后，实际求 (19) 的积分的办法是：写出该方程被 IV, p.~191 的表达式 (32) 所满足，其中 \(p_{k}\) 是次数 \(\leqslant n_{k}-1\) 、系数在 E 中的任意多项式；在如此得到的等式的两边等同 \(e^{r_{k}t}\) （对 \(1 \leqslant k \leqslant q\) ）的系数，便得到关于多项式 \(p_{k}\) 的系数的线性方程：容易证明，这些方程把次数 \textgreater0 的项确定为 \(p_{k}\) 的常数项的函数，而该常数项本身是方程 \((A - r_{k}I)^{n_{k}}.x = 0\) 的解，该方程定义子空间 \(E_{k}\) （“待定系数法”）。

注. 当存在 E 的一个基，使 A 关于该基的矩阵具有实元素时 (cf.~IV, p.~186, 注 2)，A 的特征方程就有实系数。对 E 的每个向量 \(\mathbf{x} = (\xi_k)_{1 \leqslant k \leqslant n}\) （在所考虑的基下表示），令 \(\overline{x} = (\overline{\xi}_k)_{1 \leqslant k \leqslant n}\) ；映射 \(x \mapsto \overline{x}\) 是 E 的一个反线性对合。我们知道 (Alg., VII)，若 \(r_k\) 是特征方程的一个非实根，\(E_k\) 是相应的稳定子空间，则 \(\overline{r}_k\) 是一个特征根，其重数 \(n_k\) 与 \(r_k\) 的相同，且 \(E'_k\) （即 \(E_k\) 在映射 \(x \mapsto \overline{x}\) 下的像）是对应于 \(\overline{r}_k\) 的稳定子空间。由此推出：若 \(u_j\) ( \(1 \leqslant j \leqslant n_k\) ) 是 \(n_k\) 个取值于 \(E_k\) 的线性无关的积分，则 \(2n_k\) 个积分 \(u_j + \overline{u}_j\) 、\(i(\mathbf{u}_j - \overline{\mathbf{u}}_j)\) 线性无关，并且它们关于 E 的所选基的分量都是实函数。若 \(r_k\) 是一个实特征根，则 IV, p.~186 的注 2 表明（用同样的记号），存在 \(n_k\) 个线性无关的积分 \(v_j\) ( \(1 \leqslant j \leqslant n_k\) )，它们取值于 \(E_k\) ，其分量为实值。这样便得到 (19) 的一个所有分量均为实值的基本积分组。

\subsection{n 阶线性方程}

所谓 n 阶线性微分方程，是指形如

\[
\mathrm{D} ^ {n} x - a _ {1} (t) \mathrm{D} ^ {n - 1} x - \dots - a _ {n - 1} (t) \mathrm{D} x - a _ {n} (t) x = b (t)\tag{33}
\]

的方程，其中 \(a_{k}\) ( \(1 \leqslant k \leqslant n\) ) 与 b 是实变量 t 的实（复）值函数，定义在 R 的一个区间 J 上。IV, p.~164 的一般方法表明，这个方程等价于由 n 个一阶方程组成的线性方程组

\[
\left\{ \begin{array}{l l} \frac {d x _ {k}}{d t} = x _ {k + 1} & (1 \leqslant k \leqslant n - 1) \\ \frac {d x _ {n}}{d t} = a _ {1} (t) x _ {n} + a _ {2} (t) x _ {n - 1} + \dots + a _ {n} (t) x _ {1} + b (t) \end{array} \right.\tag{34}
\]

也就是说，等价于线性方程

\[
\frac {d \mathbf {x}}{d t} = A (t). \mathbf {x} + \mathbf {b} (t)\tag{35}
\]

其中我们令 \(\mathbf{x}=(x_{1},x_{2},\ldots,x_{n})\in\mathbf{C}^{n},\mathbf{b}(t)=(0,0,\ldots,0,b(t))\) ，而矩阵 \(A(t)\) 由下式定义

\[
A (t) = \left( \begin{array}{c c c c c} 0 & 1 & 0 & \ldots & 0 \\ 0 & 0 & 1 & \ldots & 0 \\ \vdots & \vdots & \vdots & \ldots & \vdots \\ 0 & 0 & 0 & \ldots & 1 \\ a _ {n} (t) & a _ {n - 1} (t) & a _ {n - 2} (t) & \ldots & a _ {1} (t) \end{array} \right).
\]

因此，对 n 阶线性方程的研究就是把上面的一般结果应用于特殊的线性方程 (35)。对使函数 \(a_{j}\) ( \(1 \leqslant j \leqslant n\) ) 与 b 均为规则函数的每个区间 J，存在唯一的函数 u，它定义在 J 上，具有 n - 1 阶连续导数，并且在这个区间上（一个可数集的点除外）具有 n 阶规则导数，在 J 的一个可数子集的余集上满足 (33)，并且使得

\[
u (t _ {0}) = x _ {0}, \qquad \mathrm{D} u (t _ {0}) = x _ {0} ^ {\prime}, \dots , D ^ {n - 1} u (t _ {0}) = x _ {0} ^ {(n - 1)}\tag{36}
\]

其中 \(t_0\) 是 \(J\) 的任意一点，而 \(x_0, x_0', \ldots, x_0^{(n-1)}\) 是 \(n\) 个任意的复数。

为使齐次方程的 \(p\) 个积分 \(u_{j}\) ( \(1 \leqslant j \leqslant p\) )

\[
\mathrm{D} ^ {n} x - a _ {1} (t) \mathrm{D} ^ {n - 1} x - \dots - a _ {n - 1} (t) \mathrm{D} x - a _ {n} (t) x = 0\tag{37}
\]

（这个方程与 (33) 相伴）线性无关（在由 J 到 C 的连续映射组成、视为 C 上向量空间的空间 \(\mathcal{C}(\mathrm{J};\mathbf{C})\) 中），必须且只需相应的 p 个积分 \(\mathbf{u}_{j}=(u_{j},\mathrm{D}u_{j},\ldots,\mathrm{D}^{n-1}u_{j})\) （齐次方程 \(d\mathbf{x}/dt=A(t)\mathbf{x}\) 的积分）线性无关（在空间 \(\mathcal{C}(\mathrm{J};\mathbf{C}^{n})\) ——由 J 到 \(C^{n}\) 的连续映射所组成的空间——中）。显然这个条件是必要的。

反之，若存在 \(n\) 个不全为零的复常数 \(\lambda_j\) ，使得 \(\sum_{j=1}^{n} \lambda_j u_j(t) = 0\) 在 \(J\) 上恒成立，则可推出 \(\sum_{j=1}^{n} \lambda_j D^k u_j(t) = 0\) 在 \(J\) 上对每个整数 \(k\) （满足 \(1 \leqslant k \leqslant n - 1\) ）成立，这就意味着 \(\sum_{j=1}^{n} \lambda_j \mathbf{u}_j(t) = 0\) 在 \(J\) 上成立。

从而 (IV, p.~180, cor. 1)

命题 8. 定义在 J 上的齐次线性方程 (37) 的积分所成之集是域 C 上一个 n 维向量空间。

给定方程 (37) 的任意 n 个积分 \(u_{j}\) ( \(1 \leqslant j \leqslant n\) )，我们把这组积分的朗斯基行列式称为（关于 \(C^{n}\) 的典范基的）相应的 n 个积分 \(u_{j}\) （方程 \(d\mathbf{x}/dt = A(t)\mathbf{x}\) 的积分）的行列式，也就是说，函数

\[
\mathrm{W} (t) = \left| \begin{array}{c c c c} u _ {1} (t) & u _ {2} (t) & \ldots & u _ {n} (t) \\ \mathrm{D} u _ {1} (t) & \mathrm{D} u _ {2} (t) & \ldots & \mathrm{D} u _ {n} (t) \\ \vdots & \vdots & \ldots & \vdots \\ \mathrm{D} ^ {n - 1} u _ {1} (t) & \mathrm{D} ^ {n - 1} u _ {2} (t) & \ldots & \mathrm{D} ^ {n - 1} u _ {n} (t) \end{array} \right|.
\]

为使 n 个积分 \(u_{j}\) 线性无关，必须且只需 \(\mathrm{W}(t) \neq 0\) 在 J 上成立；此外，只需 \(\mathrm{W}(t_{0}) \neq 0\) 在 J 的仅一点 \(t_{0}\) 处成立即可 (IV, p.~184, prop. 4)；进一步，有 (IV, p.~184, prop. 5)

\[
\mathrm{W} (t) = \mathrm{W} (t _ {0}) \exp \left(\int_ {t _ {0}} ^ {t} a _ {1} (s) d s\right).\tag{38}
\]

我们把方程 (35) 的预解 \(C(t, t_{0})\) 与它关于 \(C^{n}\) 的典范基的矩阵等同起来；这个矩阵的各列 \(\mathbf{v}_{j}(t, t_{0})\) ( \(1 \leqslant j \leqslant n\) ) 便是 n 个线性无关的积分

\[
\mathbf {v} _ {j} (t, t _ {0}) = (v _ {j} (t, t _ {0}), \mathrm{D} v _ {j} (t, t _ {0}), \dots , \mathrm{D} ^ {n - 1} v _ {j} (t, t _ {0}))
\]

它们是齐次方程 \(d\mathbf{x}/dt = A(t)\) . x 的积分，对应于方程 (37) 的 n 个线性无关的积分 \(v_{j}(t, t_{0})\) ，这些积分使得

\[
\mathbf {D} ^ {k - 1} v _ {j} (t _ {0}, t _ {0}) = \delta_ {j k}
\]

（克罗内克记号）对 \(1 \leqslant j \leqslant n\) , \(1 \leqslant k \leqslant n\) 成立（约定 \(D^{0}v_{j} = v_{j}\) ）。特别地，由此得出，把常数变易法 (IV, p.~182) 应用于方程 (35)，这里给出 (33) 的一个特积分，它在点 \(t_{0}\) 处连同其前 n - 1 阶导数都等于 0，即函数

\[
w (t) = \int_ {t _ {0}} ^ {t} v _ {n} (t, s) b (s) d s.\tag{39}
\]

在方程 \(D^{n}x = b(t)\) 的特殊情形，公式 (39) 重新给出 \(n^{th}\) 原函数（规则函数 \(b(t)\) 的在点 \(t_{0}\) 处连同其前 n - 1 阶导数都为零者）的表达式，即

\[
w (t) = \int_ {t _ {0}} ^ {t} b (s) \frac {(t - s) ^ {n - 1}}{(n - 1) !} d s
\]

(II, p.~62, 公式 (19))：\(D^{n}x = 0\) 的在点 \(t_{0}\) 处连同其前 n - 2 阶导数都为零、且其 \(n - 1^{th}\) 阶导数在该点等于 1 的积分，实际上就是多项式 \((t - t_{0})^{n-1}/(n - 1)!\) 。

\subsection{常系数 n 阶线性方程}

若方程 (33) 中的系数 \(a_{j}\) 为常数，则相应的矩阵 \(A\) 也是常数；写出 \(e^{rt}\) 是一个解，便得到特征方程，即

\[
r ^ {n} - a _ {1} r ^ {n - 1} - \ldots - a _ {n - 1} r - a _ {n} = 0.\tag{40}
\]

设 \(r_j\) ( \(1 \leqslant j \leqslant q\) ) 是这个方程的互异根，\(n_j\) ( \(1 \leqslant j \leqslant q\) ) 是根 \(r_j \left( \sum_{j=1}^{q} n_j = n \right)\) 的重数。由 IV, p.~188 至 194 的结果，对齐次方程而言，每个根 \(r_j\) 都对应着

\[
\mathrm{D} ^ {n} x - a _ {1} \mathrm{D} ^ {n - 1} x - \dots - a _ {n - 1} \mathrm{D} x - a _ {n} x = 0\tag{41}
\]

一个由 \(n_j\) 个线性无关的积分所成的组

\[
u _ {j k} (t) = e ^ {r _ {j} t} p _ {j k} (t),
\]

其中 \(p_{jk}\) 是次数 \(\leqslant n_{j}-1\) 的（复系数）多项式（ \(1\leqslant k\leqslant n_{j}\) ）；此外，如此得到的 n 个积分 \(u_{jk}\) ( \(1\leqslant j\leqslant q\) , \(1\leqslant k\leqslant n_{j}\) ) 线性无关。由此推出，\(n_{j}\) 个多项式 \(p_{jk}\) ( \(1\leqslant k\leqslant n_{j}\) ) 在 t 的次数 \(\leqslant n_{j}-1\) 的多项式所成的空间中线性无关，从而构成该空间（在 C 上）的一个基，因为后者的维数是 \(n_{j}\) 。换言之：

命题 9. 设 \(r_j\) ( \(1 \leqslant j \leqslant q\) ) 是特征方程 (40) 的互异根，\(n_j\) 是根 \(r_j\) ( \(1 \leqslant j \leqslant q\) ) 的重数。那么 \(n\) 个函数 \(t^k e^{r_j t}\) ( \(1 \leqslant k \leqslant n_j\) , \(1 \leqslant j \leqslant q\) ) 是齐次方程 (41) 的线性无关的积分。

可以直接地证明这个结果如下。由方程 (41) 可知，这个方程的每个积分的 \(n^{th}\) 导数在 R 上可导，由此立即推出（对整数 m \textgreater{} n 用归纳法），(41) 的每个积分都有 m 阶导数，也就是说，在 R 上无穷次可导。设 D 是 R 上无穷次可导的复值函数在 C 上所成的（非拓扑）向量空间；映射 \(x \mapsto Dx\) 是这个空间的一个自同态，而方程 (41) 可以写成

\[
f (\mathrm{D}) x = 0\tag{42}
\]

\[
\text { where } f (\mathrm{D}) = \mathrm{D} ^ {n} - a _ {1} \mathrm{D} ^ {n - 1} - \dots - a _ {n - 1} \mathrm{D} - a _ {n} (\mathrm{Alg}., \mathrm{IV}, \mathrm{p}. 4).
\]

命题 10. 设 \(g\) 与 \(h\) 是两个互素多项式，使得 \(f = gh\) 。(42) 的解所成的子空间是下列两个方程的解子空间的直和：

\[
g (\mathrm{D}) x = 0, \quad h (\mathrm{D}) x = 0.
\]

现在，由贝祖恒等式 (Alg., VII. 2, th. 1)，存在两个多项式 \(p(\mathrm{D})\) 与 \(q(\mathrm{D})\) ，使得 \(p(\mathrm{D})g(\mathrm{D}) + q(\mathrm{D})h(\mathrm{D}) = 1\) 。对 (42) 的每个解 \(x\) ，可以写出 \(x = y + z\) ，其中 \(y = p(\mathrm{D})g(\mathrm{D})x\) 且 \(z = q(\mathrm{D})h(\mathrm{D})x\) ；于是 \(h(\mathrm{D})y = p(\mathrm{D})(f(\mathrm{D})x) = 0\) 且 \(g(\mathrm{D})z = q(\mathrm{D})(f(\mathrm{D})x) = 0\) 。另一方面，若 \(g(\mathrm{D})x = 0\) 与 \(h(\mathrm{D})x = 0\) 同时成立，则可推出

\[
x = p (\mathrm{D}) (g (\mathrm{D}) x) + q (\mathrm{D}) (h (\mathrm{D}) x) = 0,
\]

这就完成了证明。

用前面的记号，可以写出

\[
f (\mathrm{D}) = \prod_ {j = 1} ^ {q} (\mathrm{D} - r _ {j}) ^ {n _ {j}}
\]

而命题 10（对 \(q\) 用归纳法）表明，(42) 的解子空间是下列 \(q\) 个方程

\[
(\mathrm{D} - r _ {j}) ^ {n _ {j}} x = 0 \quad (1 \leqslant j \leqslant q).\tag{43}
\]

现在，对每个复数 r 有

\[
\mathrm{D} (e ^ {r t} x) = e ^ {r t} (\mathrm{D} + r) x\tag{44}
\]

所以 (43) 等价于

\[
\mathrm{D} ^ {n _ {j}} \left(e ^ {- r _ {j} t} x\right) = 0
\]

从而它以函数 \(e^{r_{j}t} p_{j}(t)\) 为解，其中 \(p_{j}\) 遍取次数 \(\leqslant n_{j}-1\) 的多项式所成之集；这样就重新得到 IV, p.~194 的命题 9。

假设齐次方程 (41) 已经解出（也就是说，假设特征根已经求出），我们知道常数变易法可以求出非齐次方程

\[
\mathrm{D} ^ {n} x - a _ {1} \mathrm{D} ^ {n - 1} x - \dots - a _ {n - 1} \mathrm{D} x - a _ {n} x = b (t)\tag{45}
\]

的解，其中 \(b(t)\) 是任意的规则函数 (IV, p.~182)；我们指出，若 \(b(t)\) 在一个区间 \(J\) 上无穷次可导，则 (45) 的所有积分都在 J 上无穷次可导。在特殊情形 \(b(t) = e^{\alpha t} p(t)\) （其中 p 是一个（复系数）多项式，\(\alpha\) 是一个任意复数）下，用下面的方法可以更简单地得到 (45) 的一个积分。令 \(x = e^{\alpha t} y\) ；方程

\[
f (\mathrm{D}) x = e ^ {\alpha t} p (t)
\]

由 (44) 可以写成

\[
f (\alpha + \mathrm{D}) y = p (t)
\]

或者，对多项式 \(f(\mathbf{D})\) 应用泰勒公式，

\[
\frac {f ^ {(n)} (\alpha)}{n !} \mathrm{D} ^ {n} y + \frac {f ^ {(n - 1)} (\alpha)}{(n - 1) !} \mathrm{D} ^ {n - 1} y + \dots + \frac {f ^ {\prime} (\alpha)}{1 !} \mathrm{D} y + f (\alpha) y = p (t).\tag{46}
\]

设 \(m\) 是多项式 \(p(t) = \sum_{k=0}^{m} \lambda_k t^{m-k}\) 的次数；若 \(f(\alpha) \neq 0\) （即 \(\alpha\) 不是特征根），则存在唯一的一个多项式 \(u(t) = \sum_{k=0}^{m} c_k t^{m-k}\) （次数为 \(m\) ）是 (46) 的解，因为系数 \(c_k\) 由线性方程组

\[
\begin{array}{r l} f (\alpha) c _ {k} + \binom {m - k + 1} {1} f ^ {\prime} (\alpha) c _ {k - 1} + \binom {m - k + 2} {2} f ^ {\prime \prime} (\alpha) c _ {k - 2} + \dots \\ & + \binom {m} {k} f ^ {(k)} (\alpha) c _ {0} = \lambda_ {k} \quad (0 \leqslant k \leqslant m) \end{array}
\]

所确定，而这个方程组显然有且仅有一个解。反之，若 \(\alpha\) 是一个特征根，h 是它的重数，则前面的计算表明，存在唯一的一个 m 次多项式，使得 \(D^{h}y = v(t)\) 的每个解都是一个积分；换句话说，(46) 的每个多项式解这时都是 \(m + h\) 次的（“共振”）。

\subsection{常系数线性方程组}

沿用 \(n^{\circ}\) 8 的记号，我们更一般地考虑形如下的由 m 个微分方程组成的方程组

\[
\sum_ {k = 1} ^ {n} p _ {j k} (\mathrm{D}) x _ {k} = b _ {j} (t) \quad (1 \leqslant j \leqslant m)\tag{47}
\]

的方程组，其中未知函数 \(x_{k}\) ( \(1 \leqslant k \leqslant n\) ) 与右端 \(b_{j}\) ( \(1 \leqslant j \leqslant m\) ) 都是实变量 t 的复值函数，而 \(p_{jk}(D)\) 是微分算子 D 的常（复）系数多项式（任意次数）（ \(1 \leqslant j \leqslant m, 1 \leqslant k \leqslant n\) ）。

这样的方程组与 IV, p.~1164（公式 (5)）中所考虑的那些不同型，如下例所示：

\[
\left\{ \begin{array}{c} \mathrm{D} x _ {1} = a (t) \\ \mathrm{D} ^ {2} x _ {1} + \mathrm{D} x _ {2} + x _ {2} = b (t). \end{array} \right.\tag{48}
\]

我们将把自己限于函数 \(b_{j}(t)\) 在区间 J 上无穷次可导的情形，并且只寻求在 J 上无穷次可导的解 \((x_{k})_{1\leqslant k\leqslant n}\) 。令 \(\mathbf{b}(t)=\left(b_{1}(t),\ldots,b_{m}(t)\right)\) （这是由 J 到 \(C^{m}\) 中的一个映射）以及 \(\mathbf{x}=\left(x_{1},x_{2},\ldots,x_{n}\right)\) ，方程组 (47) 便可写成

\[
P (\mathrm{D}) \mathbf {x} = \mathbf {b} (t)\tag{49}
\]

其中 \(P(\mathrm{D})\) 是矩阵 \((p_{jk}(\mathrm{D}))\) ，有 m 行 n 列，其系数属于 D 的多项式、以 C 为系数的环 C{[}D{]}。设 \(f_{j}(\mathrm{D})\) ( \(1 \leqslant j \leqslant r \leqslant \operatorname{Min}(m, n)\) ) 是矩阵 \(P(\mathrm{D})\) 的非零相似不变量；我们知道 (Alg., VII, p.~32)，它们是一些完全确定的首一多项式，使得 \(f_{j}\) 整除 \(f_{j+1}\) （对 \(1 \leqslant j \leqslant r - 1\) ；r 是 \(P(\mathrm{D})\) 的秩）；此外，存在两个方阵 \(U(\mathrm{D})\) 与 \(V(\mathrm{D})\) ，阶数分别为 m 与 n，它们是可逆的（在阶数分别为 m 与 n、以 D 的复系数多项式环 C{[}D{]} 的元素为系数的方阵环中），并且使得矩阵 \(Q(\mathrm{D}) = (q_{jk}(\mathrm{D})) = U(\mathrm{D}) P(\mathrm{D}) V(\mathrm{D})\) 的所有元素除对角元 \(q_{jj}(\mathrm{D}) = f_{j}(\mathrm{D})\) （对 \(1 \leqslant j \leqslant r\) ）外都为零。现在令 \(\mathbf{y} = V^{-1}(\mathbf{D})\mathbf{x}\) ；方程 (49) 等价于方程

\[
U (\mathrm{D}) \big (P (\mathrm{D}) \big (V (\mathrm{D}) \mathbf {y} \big) \big) = U (\mathrm{D}) \mathbf {b}
\]

也就是

\[
Q (\mathrm{D}) \mathbf {y} = U (\mathrm{D}) \mathbf {b}\tag{50}
\]

因为 \(U(\mathbf{D})\) 可逆。现在，若 \(\mathbf{y} = (y_1, y_2, \ldots, y_n)\) ，且若

\[
U (\mathrm{D}) \mathbf {b} (t) = \left(c _ {1} (t), \dots , c _ {m} (t)\right),
\]

则方程 (50) 可以写成

\[
f _ {j} (\mathrm{D}) y _ {j} = c _ {j} (t) \quad \text {   for   } 1 \leqslant j \leqslant r\tag{51}
\]

\[
0 = c _ {j} (t) \quad \text {   for   } r + 1 \leqslant j \leqslant m.\tag{52}
\]

因此，除非条件 (52) 满足，该方程组没有任何无穷次可导的解；这时 \(y_{j}\) （下标 \(j \leqslant r\) ）的确定就化为对 r 个常系数线性微分方程 (51) 的积分；而 \(y_{j}\) （下标 \textgreater{} r 的那些）是任意的无穷次可导函数。这样确定了 (50) 的解 y 之后，就由公式 \(\mathbf{x} = V(\mathrm{D})\mathbf{y}\) 得出 (47) 的解。

注. 1) 某些多项式 \(f_{j}(\mathrm{D})\) 可以化为非零常数；这时相应的 \(y_{j}\) 就被完全确定。

\begin{enumerate}
\def\labelenumi{\arabic{enumi})}
\setcounter{enumi}{1}
\item
  当 \(b_{j}\) 全为零时，即当方程组 (47) 是齐次的时，条件 (52) 总是满足的；进一步，若 \(r = n\) ，则可看出 (47) 的解所成之集是 \(\mathbf{C}\) 上的一个向量空间，其维数等于诸 \(f_{j}(\mathrm{D})\) 的次数之和，也就是 \(\det(P(\mathrm{D}))\) 的次数。
\item
  给定多项式 \(p_{jk}(\mathbf{D})\) ，一个当右端无穷次可导（或某个阶数可导）时容许解的方程组 (47)，在右端是任意规则函数时可能不容许解：例 (48) 表明了这一点，当 \(a(t)\) 没有原函数时它没有解。我们在这里不去寻求当右端是任意规则函数时出现的补充可解性条件。
\end{enumerate}

\setcounter{exercisegroup}{0}
\refstepcounter{exercisegroup}
\section*{第 \S\arabic{exercisegroup} 节习题}
\phantomsection
\addcontentsline{toc}{section}{第 \protect\S\arabic{exercisegroup} 节习题}

\begin{enumerate}
\def\labelenumi{\arabic{enumi})}
\tightlist
\item
  \begin{enumerate}
  \def\labelenumii{\alph{enumii})}
  \tightlist
  \item
    沿用 IV, p.~165 的记号，假设 \(\mathbf{f}\) 在 \(\mathrm{J} \times \mathrm{S}\) 上一致连续。不用选择公理证明 IV, p.~166 的命题 3。（设 \(\delta\) 使得关系 \(|t_2 - t_1| \leqslant \delta\) 、\(\| \mathbf{x}_2 - \mathbf{x}_1 \| \leqslant \delta\) 蕴含 \(\| \mathbf{f}(t_2, \mathbf{x}_2) - \mathbf{f}(t_1, \mathbf{x}_1) \| \leqslant \varepsilon\) ；把 J 分成长度 \(\leqslant \inf(\delta, \delta / M)\) 的区间，并用归纳法定义近似解。）
  \end{enumerate}
\end{enumerate}

\begin{enumerate}
\def\labelenumi{\alph{enumi})}
\setcounter{enumi}{1}
\tightlist
\item
  当 E 是有限维且 f 在 \(I \times H\) 上是利普希茨的时，不用选择公理证明 IV, p.~166 的命题 3。（注意，对每个 \(\delta > 0\) ，存在 S 的有限多个点 \(x_{i}\) ，使得 S 的每个点都以 \(\leqslant \delta\) 的距离靠近某个 \(x_{i}\) ；然后照 a) 的做法进行，考虑规则函数 \(t \mapsto \mathbf{f}(t, \mathbf{x}_{i})\) ，它们的个数是有限的。）
\end{enumerate}

\begin{enumerate}
\def\labelenumi{\arabic{enumi})}
\setcounter{enumi}{1}
\item
  给定两个数 \(b > 0\) 、\(M > 0\) 以及任意一个数 \(\varepsilon > 0\) ，给出一个数量微分方程 \(x' = g(x)\) 的例子，使得 \(|g(x)| \leqslant M\) （对 \(|x| \leqslant b\) ），并且它有一个积分 \(x = u(t)\) ，该积分在区间 \(] - b/M - \varepsilon, b/M + \varepsilon[\) 上连续，但在点 \(x = \frac{b}{M} + \varepsilon\) 处没有有限极限（这样定义 \(g\) ：使所考虑的积分在 \(] - b/M - \varepsilon, b/M + \varepsilon[\) 上有连续导数，且该导数等于常数 \(M\) 于 \(] - b/M, b/M[\) 上）。
\item
  设 \(S \subset H\) 是以 \(\mathbf{x}_0\) 为中心、\(r\) 为半径的开球，\(\mathbf{f}\) 是以常数 \(k > 0\) 为利普希茨常数的、在 \(I \times S\) 上的利普希茨函数；再设 \(t \mapsto \mathbf{f}(t, \mathbf{x}_0)\) 在 \(I\) 上有界，并以 \(M_0\) 表示 \(\| \mathbf{f}(t, \mathbf{x}_0) \|\) （对 \(t \in I\) ）的上确界。证明：对所有的 \(t_0 \in I\) ，存在一个积分 \(\mathbf{u}\) ——\(\mathbf{x}' = \mathbf{f}(t, \mathbf{x})\) 的积分——它取值于 \(S\) ，等于 \(\mathbf{x}_0\) （在点 \(t_0\) 处），并且定义在 \(I\) 与区间 \(]t_0 - \lambda, t_0 + \lambda[\) 的交上，其中
\end{enumerate}

\[
\lambda = \frac {1}{k} \log \left(1 + \frac {k z}{M _ {0}}\right)
\]

（注意，有 \(\| \mathbf{u}(t) - \mathbf{x}_0\| \leqslant \mathbf{M}(t - t_0) + k\int_{t_0}^t\| \mathbf{u}(s) - \mathbf{x}_0\| ds\) （对 \(t > t_0\) ））。

* 4) a) 设 \(I = ]t_0 - a, t_0 + a[\) 是 \(\mathbf{R}\) 中的开区间，\(S\) 是以 \(\mathbf{x}_0\) 为中心、\(r\) 为半径的开球（在 \(E\) 中），\(\mathbf{f}\) 是 \(I \times S\) 上的局部利普希茨函数。设 \((s, z) \mapsto h(s, z)\) 是一个 \(\geqslant 0\) 函数，对 \(0 \leqslant s \leqslant a\) 与 \(0 \leqslant z \leqslant r\) 有定义且连续，使得对每个 \(s \in [0, a]\) ，映射 \(z \mapsto h(s, z)\) 是递增的；假设 \(\| \mathbf{f}(t, \mathbf{x}) \| \leqslant h(|t - t_0|, \| \mathbf{x} - \mathbf{x}_0 \|)\) 在 \(I \times S\) 上成立。设 \(\varphi\) 是区间 \([0, \alpha]\) （\(\alpha < a\) ）上一个规则函数的原函数，取值于 \([0, r[\) ，使得 \(\varphi(0) = 0\) 且 \(\varphi'(s) > h(s, \varphi(s))\) （在 \([0, \alpha]\) 上，除一个可数集外）。证明积分 \(\mathbf{u}\) ——\(\mathbf{x}' = \mathbf{f}(t, \mathbf{x})\) 的、等于 \(\mathbf{x}_0\) 于点 \(t_0\) 的积分——定义在 \(]t_0 - \alpha, t_0 + \alpha[\) 上，并且在这个区间上 \(\| \mathbf{u}(t) - \mathbf{x}_0 \| \leqslant \varphi(|t - t_0|)\) 。

\begin{enumerate}
\def\labelenumi{\alph{enumi})}
\setcounter{enumi}{1}
\tightlist
\item
  由 a) 推出：若 f 定义在 I × E 上，且存在一个函数 h，有定义、连续、递增且 \textgreater{} 0 于 {[}0, +∞{[} 上，使得 \(\int_{0}^{+\infty} dz / h(z) = +\infty\) ，并且 \(\|\mathbf{f}(t, \mathbf{x})\| \leqslant h(\|x\|)\) ，则 \(\mathbf{x}' = \mathbf{f}(t, \mathbf{x})\) 的每个积分都在整个 I 上有定义。
\end{enumerate}

* 5) 设 \(E\) 是实数序列 \(\mathbf{x} = (x_{n})\) （满足 \(\lim_{n\to \infty}x_n = 0\) ）所成的完备赋范空间，赋以范数 \(\| \mathbf{x}\| = \sup_n|x_n|\) 。对每个整数 \(n\geqslant 0\) ，设 \(\mathbf{e}_n\) 是除下标 \(n\) 的项等于 1 外其余各项都为零的序列；空间 \(E\) 是二维子空间 \(V_{n}\) （由 \(\mathbf{e}_n\) 与 \(\mathbf{e}_{n + 1}\) 生成）与闭子空间 \(W_{n}\) （由那些 \(\mathbf{e}_k\) ——下标异于 \(n\) 与 \(n + 1\) 者——生成）的直和。设 \(\mathbf{f}_n\) 是 \(E\) 上的连续利普希茨函数，取值于 \(E\) ，在模 \(W_{n}\) 的每个等价类上为常数，等于 \(\mathbf{e}_{n + 1} - \mathbf{e}_n\) 于连接 \(\mathbf{e}_n\) 与 \(\mathbf{e}_{n + 1}\) 的直线上，并且在 \(V_{n}\) 与球 \(\| \mathbf{x}\| \leqslant \frac{1}{4}\) 的交上等于 0。另一方面，对每个整数 \(n > 0\) ，设 \(\varphi_{n}\) 是一个实函数，在区间 \(\left[\frac{1}{n + 1},\frac{1}{n}\right]\) 上有定义且连续，在该区间的端点处等于 0，并且使得 \(\int_{1 / n + 1}^{1 / n}\varphi_n(t)dt = 1\) 。设 \(I\) 是区间 {[}0, 1{]} （在 \(\mathbf{R}\) 中）；在 \(I\times E\) 上考虑函数 \(\mathbf{f}\) ，它取值于 \(E\) ，定义如下：\(\mathbf{f}(0,\mathbf{x}) = 0\) （对每个 \(\mathbf{x}\in E\) ）；对 \(\frac{1}{n + 1}\leqslant t\leqslant \frac{1}{n}\) ，令 \(\mathbf{f}(t,\mathbf{x}) = \varphi_n(t)\mathbf{f}_n(\mathbf{x})\) 。证明 \(\mathbf{f}\) 在 \(I\times E\) 上连续且局部利普希茨，但存在一个积分 \(\mathbf{u}\) ——方程 \(\mathbf{x}' = \mathbf{f}(t,\mathbf{x})\) 的积分——它在 {[}0, 1{]} 上有定义且有界，等于 \(\mathbf{e}_n\) （当 \(t = 1 / n\) ），从而当 \(t\) 趋于 0 时没有极限。

* 6) a) 设 \(\mathrm{I} = [0, +\infty[\) ；若 \(\mathbf{f}\) 在 \(\mathrm{I} \times \mathrm{E}\) 上以一个规则函数 \(k(t) > 0\) 为利普希茨函数，且积分 \(\int_0^{+\infty} k(t) dt\) 收敛，证明方程 \(\mathbf{x}' = \mathbf{f}(t, \mathbf{x})\) 的每个积分都在整个 \(\mathrm{I}\) 上有定义。

\begin{enumerate}
\def\labelenumi{\alph{enumi})}
\setcounter{enumi}{1}
\tightlist
\item
  进一步，若积分 \(\int_{0}^{+\infty}\|\mathbf{f}(t,\mathbf{x}_{0})\|dt\) （对某一点 \(x_{0}\in E\) ）收敛，证明 \(\mathbf{x}^{\prime}=\mathbf{f}(t,\mathbf{x})\) 的每个积分当 t 趋于 \(+\infty\) 时趋于一个有限极限（先证明每个积分当 t 趋于 \(+\infty\) 时有界）。
\end{enumerate}

\begin{enumerate}
\def\labelenumi{\arabic{enumi})}
\setcounter{enumi}{6}
\tightlist
\item
  考虑数量微分方程组
\end{enumerate}

\[
\frac {d x _ {i}}{d t} = \sum_ {j, k = 1} ^ {n} c _ {i j k} x _ {j} x _ {k} \quad (1 \leqslant i \leqslant n)
\]

其中 \(c_{ijk}\) 是满足 \(c_{kji} = -c_{ijk}\) 的常数。证明这个方程组的积分对 t 的一切值都有定义（注意 \(\sum_{i=1}^{n} x_{i}^{2}\) 对每个积分 \(\mathbf{x} = (x_{i})\) 都是常数）。

\begin{enumerate}
\def\labelenumi{\arabic{enumi})}
\setcounter{enumi}{7}
\tightlist
\item
  设 \(\mathbf{f}\) 是定义在 \(\mathrm{I} \times \mathrm{H}\) 上、满足 IV, p.~165 的引理 1 的条件的函数，并且对常数 \(k\) （ \(0 < k < 1\) ）与点 \(t_0 \in \mathrm{I}\) ，对所有的 \(t \in \mathrm{I}\) 与 \(\mathbf{x}_1\) 、\(\mathbf{x}_2 \in \mathrm{H}\) 有
\end{enumerate}

\[
\left\| \mathbf {f} (t, \mathbf {x} _ {1}) - \mathbf {f} (t, \mathbf {x} _ {2}) \right\| \leqslant \frac {k}{\left| t - t _ {0} \right|} \left\| \mathbf {x} _ {1} - \mathbf {x} _ {2} \right\|.
\]

证明：若 \(\mathbf{u}\) 与 \(\mathbf{v}\) 是方程 \(\mathbf{x}' = \mathbf{f}(t, \mathbf{x})\) 的两个近似解，精度分别为 \(\varepsilon_1\) 与 \(\varepsilon_2\) ，都定义在 I 上，取值于 H，并且在点 \(t_0\) 处取相同的值，则对每个 \(t \in I\) 有

\[
\left\| \mathbf {u} (t) - \mathbf {v} (t) \right\| \leqslant \frac {\varepsilon_ {1} - \varepsilon_ {2}}{1 - k} | t - t _ {0} |.
\]

由此推出，在所指出的条件下，定理 1 (IV, p.~171) 与定理 2 (IV, p.~172) 对方程 \(\mathbf{x}' = \mathbf{f}(t, \mathbf{x})\) 仍然成立。

\begin{enumerate}
\def\labelenumi{* \arabic{enumi})}
\setcounter{enumi}{8}
\tightlist
\item
  设 I 是 \(\mathbf{R}\) 中的一个区间，\(t_0\) 是 I 的一点，S 是 E 中半径 \(r\) 、中心 \(\mathbf{x}_0\) 的球，G 是由 \(\mathrm{I} \times \mathrm{S}\) 到 E 中的有界映射所成的赋范空间，其中一个映射 \(\mathbf{f}\) 的范数 \(\| \mathbf{f} \|\) 是 \(\| \mathbf{f}(t, \mathbf{x}) \|\) 在 \(\mathrm{I} \times \mathrm{S}\) 上的上确界。对每个 \(M > 0\) ，设 \(\mathrm{G}_{\mathrm{M}}\) 是 G 中的球 \(\| \mathbf{f} \| \leqslant \mathrm{M}\) 。设 L 是 G 中由 \(\mathrm{I} \times \mathrm{S}\) 到 E 中的利普希茨映射所成的子集；对每个函数 \(\mathbf{f} \in \mathrm{L} \cap \mathrm{G}_{\mathrm{M}}\) ，设 \(\mathbf{u} = \mathrm{U}(\mathbf{f})\) 是方程 \(\mathbf{x}' = \mathbf{f}(t, \mathbf{x})\) 的满足 \(\mathbf{u}(t_0) = \mathbf{x}_0\) 的积分，它取值于 S，且定义在交 \(\mathrm{J}_{\mathrm{M}}\) 上，这个交是 I 与区间
\end{enumerate}

\[
\left] t _ {0} - \frac {r}{\mathbf {M}}, t _ {0} + \frac {r}{\mathbf {M}} \right[ \text {(th. 1).}
\]

\begin{enumerate}
\def\labelenumi{\alph{enumi})}
\tightlist
\item
  设 \((\mathbf{f}_{n})\) 是属于 \(L \cap G_{M}\) 的函数序列；若 \(f_{n}\) 在 \(I \times S\) 上一致收敛于函数 f，则序列 \(\mathbf{u}_{n} = \mathrm{U}(\mathbf{f}_{n})\) 在由 \(J_{M}\) 到 E 中的有界映射所成的空间 F 中的每个聚点（关于紧收敛拓扑）都是 \(\mathbf{x}' = \mathbf{f}(t, \mathbf{x})\) 的取值 \(x_{0}\) 于点 \(t_{0}\) 的积分。反之，\(\mathbf{x}' = \mathbf{f}(t, \mathbf{x})\) 的每个定义在 \(J_{M}\) 上且满足 \(\mathbf{v}(t_{0}) = \mathbf{x}_{0}\) 的积分 v，也是某个方程 \(\mathbf{x}' = \mathbf{g}(t, \mathbf{x})\) 的积分，其中 g 是利普希茨的并且在 G 中任意接近于 f（考虑方程
\end{enumerate}

\[
\mathbf {x} ^ {\prime} = \mathbf {f} _ {n} (t, \mathbf {x}) + \mathbf {v} ^ {\prime} (t) - \mathbf {f} _ {n} (t, \mathbf {v} (t))).
\]

\begin{enumerate}
\def\labelenumi{\alph{enumi})}
\setcounter{enumi}{1}
\tightlist
\item
  进一步假设 E 是有限维的。证明：若 \(f \in G_{M}\) 满足引理 1 的假设，则对每个 \(t \in J_{M}\) ，\(H(t)\) ——即 \(\mathbf{x}' = \mathbf{f}(t, \mathbf{x})\) 的取值 \(x_{0}\) 于点 \(t_{0}\) 的那些积分在点 t 处的值所成之集——是紧连通集（要看出 \(H(t)\) 是闭的，用阿斯科利定理；要看出它是连通的，用 a)：若 \(x_{1}, x_{2}\) 是 \(H(t)\) 的两点且 \(\varepsilon > 0\) 任意，证明存在一个连通集 \(\Phi\) （由属于 \(L \cap G_{M}\) 的函数 g 组成），使得 \(\|f - g\| \leqslant \varepsilon\) 对每个函数 \(g \in \Phi\) 成立，并且诸函数 \(U(\mathbf{g})\) 在点 t 处的值所成之集是连通的且包含 \(x_{1}\) 与 \(x_{2}\) 。最后沿一个比 \(R_{+}\) 中 0 的邻域滤子更细的超滤子过渡到极限来完成证明）。
\end{enumerate}

\begin{enumerate}
\def\labelenumi{\arabic{enumi})}
\setcounter{enumi}{9}
\tightlist
\item
  设 \(f\) 是定义在长方体 \(\mathrm{P}: |t - t_0| < a\) , \(|x - x_0| < b\) （属于 \(\mathbf{R}^2\) ）上的有界连续实函数。设 \(\mathbf{M}\) 是 \(|f(t,x)|\) 在 \(\mathrm{P}\) 上的上确界，并设 \(\mathrm{I} = ]t_0 - \alpha, t_0 + \alpha[\) ，其中 \(\alpha = \inf(a, b/\mathrm{M})\) 。证明：积分集 \(\Phi\) ——即 \(x' = f(t,x)\) 的、定义在 \(\mathrm{I}\) 上且取值 \(x_0\) 于点 \(t_0\) 的积分所成之集——的上包络与下包络仍是 \(x' = f(t,x)\) 的积分，分别称为这个方程对应于点 \((t_0, x_0)\) 的最大积分与最小积分（注意，集 \(\Phi\) 对 \(\mathrm{I}\) 上的一致收敛拓扑是等度连续且闭的）。
\end{enumerate}

对每个 \(\tau \in I\) ，设 \(\xi\) 是最小积分（对应于点 \((t_0, x_0)\) 的）在点 \(\tau\) 处的值。证明：对应于点 \((\tau, \xi)\) 的最小积分与对应于点 \((t_0, x_0)\) 的最小积分在一个形如 \([\tau, \tau + h[\) 的区间上（若 \(\tau > t_0\) ）、或在一个形如 \(]\tau - h, \tau]\) 的区间上（若 \(\tau < t_0\) ）恒等。

由此推出：存在一个最大的开区间 \(]t_{1}, t_{2}[\) ，它包含 \(t_{0}\) 且含于 \(]t_{0}-a, t_{0}+a[\) ，使得对应于 \((t_{0}, x_{0})\) 的最小积分 u 可以连续地延拓到 \(]t_{1}, t_{2}[\) 上，且在 \(]t_{1}, t_{2}[\) 的每一点 t 处，值 \(u(t)\) 属于 \(]x_{0}-b, x_{0}+b[\) ，并且 u 与对应于点 \((t, u(t))\) 的最小积分在一个形如 \([t, t+h[\) 的区间上（若 \(t > t_{0}\) ）、或在一个形如 \(]t-h, t]\) 的区间上（若 \(t < t_{0}\) ）恒等；进一步证明：或者 \(t_{1} = t_{0} - a\) （相应地 \(t_{2} = t_{0} + a\) ），或者 \(\lim_{t \to t_{1}} u(t) = x_{0} \pm b\) （相应地 \(\lim_{t \to t_{2}} u(t) = x_{0} \pm b\) ）。

\begin{enumerate}
\def\labelenumi{\arabic{enumi})}
\setcounter{enumi}{10}
\tightlist
\item
  \begin{enumerate}
  \def\labelenumii{\alph{enumii})}
  \tightlist
  \item
    在长方体 P: \(|t - t_{0}| < a\) , \(|x - x_{0}| < b\) 上，设 g 与 h 是两个连续实函数，使得在 P 上 \(g(t, x) < h(t, x)\) 。设 u（相应地 v）是 \(x' = g(t, x)\) （相应地 \(x' = h(t, x)\) ）的满足 \(u(t_{0}) = x_{0}\) （相应地 \(v(t_{0}) = x_{0}\) ）的积分，定义在区间 \([t_{0}, t_{0} + c[;\) 上；证明对 \(t_{0} < t < t_{0} + c\) 有 \(u(t) < v(t)\) （考虑使这个不等式成立的 t 所成之集的上确界 \(\tau\) ）。
  \end{enumerate}
\end{enumerate}

\begin{enumerate}
\def\labelenumi{\alph{enumi})}
\setcounter{enumi}{1}
\item
  设 u 是 \(x' = g(t, x)\) 的对应于点 \((t_0, x_0)\) 的最大积分 (exerc. 10)，定义在区间 \([t_0, t_0 + c[\) 上，取值于 \(|x - x_0| < b\) 。证明：在每个紧区间 \([t_0, t_0 + d]\) （含于 \([t_0, t_0 + c[\) ）中，方程 \(x' = g(t, x) + \varepsilon\) 的对应于点 \((t_0, x_0)\) 的最小积分与最大积分只要 \(\varepsilon > 0\) 足够小就有定义，并且当 \(\varepsilon\) 沿 \textgreater{} 0 的值趋于 0 时一致收敛于 u。
\item
  设 g 与 h 是定义在 P 上的两个连续实函数，且在 P 上 \(g(t, x) \leqslant h(t, x)\) 。设 \([t_0, t_0 + c[\) 是一个区间，在其上定义了 \(x' = g(t, x)\) 的一个满足 \(u(t_0) = x_0\) 的积分 u，以及 \(x' = h(t, x)\) 的对应于点 \((t_0, x_0)\) 的最大积分 v。证明在这个区间上 \(u(t) \leqslant v(t)\) （利用 b) 化为情形 a)）。
\end{enumerate}

\begin{enumerate}
\def\labelenumi{\arabic{enumi})}
\setcounter{enumi}{11}
\tightlist
\item
  设 \(u\) 是方程 \(x' = \lambda + \frac{x^2}{1 + t^2}\) 的当 \(t = 0\) 时等于 0 的积分，并设 \(J\) 是以 0 为左端点、\(u\) 在其上连续的最大区间。
\end{enumerate}

\begin{enumerate}
\def\labelenumi{\alph{enumi})}
\item
  证明：若 \(\lambda \leqslant \frac{1}{4}\) ，则 \(J = [0, +\infty[\) （用 IV, p.~199 的 exerc. 4）。
\item
  证明：若 \(\lambda >\frac{1}{4}\) ，则 \(J = [0,a[\) ，其中
\end{enumerate}

\[
\sinh \frac {\pi}{2 \sqrt {\lambda}} <   a <   \sinh \frac {\pi}{\sqrt {4 \lambda - 1}}
\]

（令 \(x = y\sqrt{1 + t^{2}}\) 并用 exerc. 11）。

* 13) a) 设 \(\mathrm{I} = [t_0, t_0 + c]\) ，并设 \(\omega\) 是一个连续实函数，\(\geqslant 0\) ，定义在 \(\mathrm{I} \times \mathbf{R}\) 上。设 \(\mathrm{S}\) 是一个以 \(\mathbf{x}_0\) 为中心的球，位于一个完备赋范空间 \(\mathrm{E}\) 中，并设 \(\mathbf{f}\) 是一个由 \(\mathrm{I} \times \mathrm{S}\) 到 \(\mathrm{E}\) 中的连续映射，使得对任意 \(t \in \mathrm{I}\) 、\(\mathbf{x}_1 \in \mathrm{S}\) 与 \(\mathbf{x}_2 \in \mathrm{S}\) 有 \(\| \mathbf{f}(t, \mathbf{x}_1) - \mathbf{f}(t, \mathbf{x}_2) \| \leqslant \omega(t, \| \mathbf{x}_1 - \mathbf{x}_2 \|)\) 。设 \(\mathbf{u}\) 与 \(\mathbf{v}\) 是 \(\mathbf{x}' = \mathbf{f}(t, \mathbf{x})\) 的两个积分，定义在 \(\mathrm{I}\) 上，取值于 \(\mathrm{S}\) ，且满足 \(\mathbf{u}(t_0) = \mathbf{x}_1\) 、\(\mathbf{v}(t_0) = \mathbf{x}_2\) ；设 \(w\) 是 \(z' = \omega(t, z)\) 的对应于点 ( \(t_0, \| \mathbf{x}_1 - \mathbf{x}_2 \|\) ) 的最大积分 (exerc. 10)，假设它定义在 \(\mathrm{I}\) 上；证明在 \(\mathrm{I}\) 上有 \(\| \mathbf{u}(t) - \mathbf{v}(t) \| \leqslant w(t)\) 。（设 \(w(t, \varepsilon)\) 是 \(z' = \omega(t, z) + \varepsilon\) 的对应于点 ( \(t_0, \| \mathbf{x}_1 - \mathbf{x}_2 \|\) ) 的最大积分，其中 \(\varepsilon > 0\) 足够小。证明 \(\| \mathbf{u}(t) - \mathbf{v}(t) \| \leqslant w(t, \varepsilon)\) 对每个 \(\varepsilon > 0\) 成立（用反证法）。）

\begin{enumerate}
\def\labelenumi{\alph{enumi})}
\setcounter{enumi}{1}
\tightlist
\item
  在 a) 的假设陈述中，把 I 换成区间 \(I' = ]t_{0} - c, t_{0}]\) 。证明：若 w 是 \(z' = \omega(t, z)\) 的对应于点 \((t_{0}, \| \mathbf{x}_{1} - \mathbf{x}_{2}\|)\) 的在这个区间上的最小积分，则 \(\| \mathbf{u}(t) - \mathbf{v}(t) \| \geqslant w(t)\) 在 \(I'\) 上成立（同样方法）。
\end{enumerate}

* 14) a) 设 \(\omega(t,z)\) 是一个连续实函数，\(\geqslant 0\) ，对 \(0 < t \leqslant a\) 与 \(z \geqslant 0\) 有定义。假设 \(z = 0\) 是 \(z' = \omega(t,z)\) 的唯一积分，它对 \(0 < t \leqslant a\) 有定义，且满足 \(\lim_{t \to 0} z(t) = 0\) 与 \(\lim_{t \to 0} z(t)/t = 0\) 。设 \(\mathrm{I} = [\mathrm{t}_0, t_0 + a[\) ，设 \(\mathrm{S}\) 是一个以 \(\mathbf{x}_0\) 为中心的球（在 \(\mathrm{E}\) 中），\(\mathbf{f}\) 是一个由 \(\mathrm{I} \times \mathrm{S}\) 到 \(\mathrm{E}\) 中的映射，使得对任意 \(t \in \mathrm{I}\) 、\(\mathbf{x}_1 \in \mathrm{S}\) 与 \(\mathbf{x}_2 \in \mathrm{S}\) 有 \(\|\mathbf{f}(t, \mathbf{x}_1) - \mathbf{f}(t, \mathbf{x}_2)\| \leqslant \omega(|t - t_0|, \| \mathbf{x}_1 - \mathbf{x}_2 \|)\) 。证明：在一个以 \(t_0\) 为左端点、含于 \(\mathrm{I}\) 的区间上，方程 \(\mathbf{x}' = \mathbf{f}(t, \mathbf{x})\) 只能有一个解 \(\mathbf{u}\) 满足 \(\mathbf{u}(t_0) = \mathbf{x}_0\) 。（用反证法：若 \(\mathbf{v}\) 是 \(\mathbf{x}' = \mathbf{f}(t, \mathbf{x})\) 的第二个积分，从下方估计 \(\|\mathbf{u}(t) - \mathbf{v}(t)\|\) 于一个以 \(t_0\) 为左端点的区间上（借助 exerc. 13 b)），从而得出矛盾。）

应用于 \(\omega(t, z) = k(z/t)\) （其中 \(0 \leqslant k < 1\) ）的情形 (cf.~IV, p.~200, exerc. 8)。

\begin{enumerate}
\def\labelenumi{\alph{enumi})}
\setcounter{enumi}{1}
\item
  a) 的结果适用于 \(\omega(t,z)=z/t\) ；但在这种情形用一个例子证明：若 u、v 是以 \(\varepsilon\) 为精度的 \(\mathbf{x}^{\prime}=\mathbf{f}(t,\mathbf{x})\) 的两个近似积分，取值 \(x_{0}\) 于点 \(t_{0}\) ，则不可能用一个只依赖于 t（而不依赖于函数 f）的数从上方控制 \(\|\mathbf{u}(t)-\mathbf{v}(t)\|\) 。（取 f 为一个由 \(R_{+}\times R\) 到 R 中的连续映射，它对 \(t\geqslant\alpha\) 以及对 \(0<t<\alpha\) 且 \(|x|\leqslant t^{2}/(\alpha-t)\) 等于 x/t，而在其余的点 \((t,x)\) 处与 x 无关。）
\item
  设 \(\theta\) 是一个连续实函数，\(\geqslant 0\) 于区间 \([0, a]\) 上。证明：若积分 \(\int_0^a \frac{\theta(t)}{t} dt\) 收敛，则 \(a)\) 的结果适用于 \(\omega(t, z) = \frac{1 + \theta(t)}{t} z\) ；另一方面，若 \(\int_0^a \frac{\theta(t)}{t} dt\) 为无穷，给出一个连续函数 \(\mathbf{f}\) 的例子，使得 \(\| \mathbf{f}(t, \mathbf{x}_1) - \mathbf{f}(t, \mathbf{x}_2) \| \leqslant \frac{1 + \theta(|t - t_0|)}{|t - t_0|} \| \mathbf{x}_1 - \mathbf{x}_2 \|\) ，但方程 \(\mathbf{x}' = \mathbf{f}(t, \mathbf{x})\) 有无穷多个积分等于 \(\mathbf{x}_0\) 于点 \(t_0\)（方法与 \(b\) 的类似）。
\end{enumerate}

\begin{enumerate}
\def\labelenumi{\arabic{enumi})}
\setcounter{enumi}{14}
\item
  设 \(f\) 是一个实函数，对 \(|t| \leqslant a\) , \(|x| \leqslant b\) 有定义且连续，使得 \(f(t,x) < 0\) （对 \(tx > 0\) ）且 \(f(t,x) > 0\) （对 \(tx < 0\) ）；证明 \(x = 0\) 是方程 \(x' = f(t,x)\) 的在点 \(t = 0\) 处取值 0 的唯一积分（用反证法）。
\item
  设 E 是一个有限维向量空间，f 是 I×H 上的有界函数，满足 IV, p.~165 的引理 1 的条件，使得方程 \(\mathbf{x}^{\prime}=\mathbf{f}(t,\mathbf{x})\) 有唯一的解 u，它定义在 I 上，取值于 H，等于 \(x_{0}\) 于点 \(t_{0}\) 。进一步假设：对足够大的 n，存在一个近似积分 \(u_{n}\) ——\(\mathbf{x}^{\prime}=\mathbf{f}(t,\mathbf{x})\) 的、精度为 1/n 的积分——定义在 I 上，取值于 H，等于 \(x_{0}\) 于点 \(t_{0}\) 。证明序列 \((\mathbf{u}_{n})\) 在含于 I 的每个紧区间上一致收敛于 u（利用序列 \((\mathbf{u}_{n})\) 在 I 上等度连续这一事实）。
\item
  为研究方程 \(x' = \sin tx\) (IV, p.~174, 例 4)，令 u = xy，并考虑相应的方程 \(u' = \frac{u}{t} + t \sin u = \mathrm{F}(t, u)\) 的为 \textgreater0 （对 t \textgreater0 ）的解（对每个这种解有 \(u(0) = 0\) ）。以 \(\Gamma_k\) 表示由关系 \((2k - 1)\pi < u < 2k\pi\) , \(t \sin u + \frac{u}{t} = 0\) （对每个整数 \(k \geqslant 1\) ）所定义的曲线；以 \(D_k\) 表示由 \((2k - 1)\pi < u < 2k\pi\) , \(t \sin u + \frac{u}{t} < 0\) 定义的开集；以 E 表示诸 \(\overline{D}_k\) 的并在 \((t, u)\) （满足 t \textgreater0 且 u \textgreater0）所成之集中的余集。
\end{enumerate}

\begin{enumerate}
\def\labelenumi{\alph{enumi})}
\item
  证明：每条与直线 \(u = 2k\pi\) 相交的积分曲线也与直线 \(u = (2k + 1)\pi\) 相交。
\item
  若积分曲线 C 与曲线 \(\Gamma_{k}\) 交于一点 \((t_{0}, u_{0})\) ，则函数 u 对 0 \textless{} t \textless{} t \textless{} \(t_{0}\) 递增，对 \(t > t_{0}\) 递减，且当 t 趋于 \(+\infty\) 时 \(u(t)\) 趋于 \((2k - 1)\pi\) ，而曲线 C 保持含于 \(D_{k}\) 中（对 \(t > t_{0}\) ）。
\item
  证明：不存在包含在 E 中且使得 \(u(t)\) 当 t 趋于 \(+\infty\) 时趋于 \(+\infty\) 的积分曲线 C。（建立沿 C 的 x 与 u 之间的微分方程 dx/du = G(u, x)。若 C 整个位于 E 中，则 \(x^{2} > u\) （对 \(u = (2k - \frac{1}{2})\pi\) ）；另一方面，利用前面的微分方程估计 \(x^{2}(u)\) ，并得出矛盾。）
\end{enumerate}

\(d\)) 证明：对每个整数 \(k\) ，存在一条且仅一条包含在 E 中且使得 \(u(t)\) 趋于 \(2k\pi\) （当 \(t\) 趋于 \(+\infty\) 时）的积分曲线 C。（令 \(v = 2k\pi - u\) ，得到方程 \(v' = \frac{v - 2k\pi}{t} + t \sin v\) ；把该方程与 \(v' = \frac{v - 2k\pi}{t} + tv\) 相比较，其中 \(t\) 趋于 \(+\infty\) 而 \(v\) 趋于 0。）

\begin{enumerate}
\def\labelenumi{\arabic{enumi})}
\setcounter{enumi}{17}
\tightlist
\item
  设 E 是实数序列 \(\mathbf{x}=(x_{n})_{n\in\mathbb{N}}\) （满足 \(\lim_{n\to\infty}x_{n}=0\) ）所成的空间，赋以范数 \(\|x\|=\sup_{n}|x_{n}|\) ，它实际上是一个完备赋范空间。对每个 \(\mathbf{x}=(x_{n})\in\mathrm{E}\) ，设 \(\mathbf{y}=\mathbf{f}(\mathbf{x})\) 是 E 的元素 \((y_{n})\) ，由 \(y_{n}=|x_{n}|^{\frac{1}{2}}+\frac{1}{n+1}\) 定义；函数 f 在 E 上连续。证明：微分方程 \(\mathbf{x}^{\prime}=\mathbf{f}(\mathbf{x})\) 没有定义在 R 中 0 的一个邻域上、取值于 E 且在点 t=0 处等于 0 的解 (cf.~IV, p.~202, exerc. 11)。
\end{enumerate}

\setcounter{exercisegroup}{1}
\refstepcounter{exercisegroup}
\section*{第 \S\arabic{exercisegroup} 节习题}
\phantomsection
\addcontentsline{toc}{section}{第 \protect\S\arabic{exercisegroup} 节习题}

\begin{enumerate}
\def\labelenumi{\arabic{enumi})}
\item
  设 E 是 R 上的一个完备赋范空间，F 是一个拓扑空间，J 是 R 的一个不退化为一点的区间；设 \((t, \xi) \mapsto A(t, \xi)\) 是一个由 \(J \times F\) 到 \(\mathcal{L}(E)\) 中的映射，使得当 \(\xi\) 趋于 \(\xi_{0}\) 时，\(A(t, \xi)\) 在 J 上一致趋于 \(A(t, \xi_{0})\) 。若 \(C(t, t_{0}, \xi)\) 是线性方程 \(d\mathbf{x}/dt = A(t, \xi)\) . x 的预解，证明：对每个紧区间 \(K \subset J\) ，\(C(t, t_{0}, \xi)\) 一致趋于 \(C(t, t_{0}, \xi_{0})\) 于 \(K \times K\) 上（当 \(\xi\) 趋于 \(\xi_{0}\) 时）。
\item
  设 \(t \mapsto A(t)\) 是一个由 \(J\) 到 \(\mathcal{L}(E)\) 中的规则映射，使得对任意两点 \(s, t\) （在 \(J\) 中），\(A(s)\) 与 \(A(t)\) 可交换。令 \(B(t) = \int_{t_0}^{t} A(s) ds\) 。证明预解 \(C(t, t_0)\) ——方程 \(d\mathbf{x}/dt = A(t)\mathbf{x}\) 的预解——等于 \(\exp(B(t))\) 。
\end{enumerate}

若 \(t \mapsto A_1(t)\) 、\(t \mapsto A_2(t)\) 是两个由 J 到 \(\mathcal{L}(\mathrm{E})\) 中的规则函数，使得对任意两点 \(s\) 、\(t\) （在 J 中），\(A_1(s)\) 、\(A_1(t)\) 、\(A_2(s)\) 、\(A_2(t)\) 两两可交换，证明

\[
\exp \left(\int_ {t _ {0}} ^ {t} \left(A _ {1} (s) + A _ {2} (s)\right) d s\right) = \exp \left(\int_ {t _ {0}} ^ {t} A _ {1} (s) d s\right) \exp \left(\int_ {t _ {0}} ^ {t} A _ {2} (s) d s\right).
\]

\begin{enumerate}
\def\labelenumi{\arabic{enumi})}
\setcounter{enumi}{2}
\tightlist
\item
  给定两个矩阵
\end{enumerate}

\[
A = \left( \begin{array}{c c} 0 & 1 \\ 0 & 0 \end{array} \right), \qquad \qquad B = \left( \begin{array}{c c} 0 & 0 \\ 1 & 0 \end{array} \right)
\]

证明 \(\exp(A+B)\neq\exp(A)\exp(B)\) 。

\begin{enumerate}
\def\labelenumi{\arabic{enumi})}
\setcounter{enumi}{3}
\item
  若 \(P\) 是 \(\mathcal{L}(\mathrm{E})\) 中任意一个自同构，证明 \(\exp(PAP^{-1}) = P\exp(A)P^{-1}\) 对每个 \(A \in \mathcal{L}(\mathrm{E})\) 成立。
\item
  用例子证明：线性方程 \(d\mathbf{x} / dt = A.\mathbf{x} + e^{\alpha t}\mathbf{p}(t)\) （其中 \(A \in \mathcal{L}(\mathrm{E})\) 与 \(t\) 无关，\(\mathbf{p}\) 是一个多项式（系数在 E 中））可以有一个积分等于一个与 \(\mathbf{p}\) 同次数的多项式，即使当 \(\alpha\) 是 \(A\) 的特征根时也是如此。
\item
  设 \(f(X)\) 是一个 n 次复系数多项式，并设
\end{enumerate}

\[
\frac {1}{f (\mathrm{X})} = \sum_ {j = 1} ^ {q} \sum_ {h = 1} ^ {n _ {j}} \frac {\lambda_ {j h}}{(\mathrm{X} - r _ {j}) ^ {h}}
\]

是有理分式 \(1 / f(\mathrm{X})\) 的部分分式分解 (Alg., VII. 7, n° 2)。证明：对任意规则函数 \(b\) ，函数

\[
t \mapsto \sum_ {j = 1} ^ {q} \sum_ {h = 1} ^ {n _ {j}} \lambda_ {j h} \int_ {t _ {0}} ^ {t} \frac {(t - s) ^ {h - 1}}{(h - 1) !} e ^ {r _ {j} (t - s)} b (s) d s
\]

是方程 \(f(\mathrm{D})x = b(t)\) 的一个积分。

\begin{enumerate}
\def\labelenumi{\arabic{enumi})}
\setcounter{enumi}{6}
\tightlist
\item
  \begin{enumerate}
  \def\labelenumii{\alph{enumii})}
  \tightlist
  \item
    设 \(t \mapsto A(t)\) 是一个由区间 \(J \subset R\) 到 \(\mathcal{L}(E)\) 中的映射，使得对每个 \(x \in E\) ，映射 \(t \mapsto A(t).x\) （由 J 到 E 中）连续。证明 \(t \mapsto \|A(t)\|\) 在每个紧区间 \(K \subset J\) 上有界（用 Gen.~Top., IX, p.~194, th. 2）。在这些条件下证明：对每一点 \((t_{0}, \mathbf{x}_{0}) \in \mathbf{J} \times \mathbf{E}\) ，方程 \(d\mathbf{x}/dt = A(t).\mathbf{x} + \mathbf{b}(t)\) 有一个且仅有一个解，它定义在 J 上且等于 \(x_{0}\) 于点 \(t_{0}\) 。进一步，若以 \(\mathbf{u}(t, t_{0}, \mathbf{x}_{0})\) 表示这个解，则映射 \(\mathbf{x}_{0} \mapsto \mathbf{u}(t, t_{0}, \mathbf{x}_{0})\) 是 E 到其自身上的一个双射双连续线性映射 \(C(t, t_{0})\) ，它满足关系式 (6) (IV, p.~179) 与 (7) (IV, p.~181)；进一步，映射 \((s, t) \mapsto C(s, t)\) 是由 \(J \times J\) 到 \(\mathcal{L}(E)\) 中的连续映射。
  \end{enumerate}
\end{enumerate}

\begin{enumerate}
\def\labelenumi{\alph{enumi})}
\setcounter{enumi}{1}
\tightlist
\item
  取 E 为实数序列 \(\mathbf{x}=(x_{n})_{n\in\mathbf{N}}\) （满足 \(\lim_{n\to\infty}x_{n}=0\) ）所成的空间，赋以范数 \(\|x\|=\sup_{n}|x_{n}|\) ，E 在其下是完备的。对每个 \(t\in J=[0,1]\) ，设 \(A(t)\) 是 E 到其自身的、使得
\end{enumerate}

\[
A (t). \mathbf {x} = \left(\frac {1}{1 + n t} x _ {n}\right) _ {n \in \mathbf {N}}.
\]

证明 \(A(t)\) 满足 a) 的条件，但映射 \(t \mapsto A(t)\) （由 J 到 \(\mathcal{L}(\mathrm{E})\) 中）在点 t = 0 处不连续，并且预解 \(C(t, t_0)\) ——方程 \(d\mathbf{x}/dt = A(t)\) . x 的预解——使得映射 \(t \mapsto C(t, t_0)\) （由 J 到 \(\mathcal{L}(\mathrm{E})\) 中）在点 t = 0 处不可导。

* 8) 设 G 是域 R 上的一个有单位元 e 的完备赋范代数。

\begin{enumerate}
\def\labelenumi{\alph{enumi})}
\item
  设 \(t \mapsto \mathbf{a}(t)\) 是区间 \(J \subset R\) 上的一个取值于 G 的规则函数。证明线性方程 \(d\mathbf{x}/dt = \mathbf{a}(t)\mathbf{x}\) 的取值 e 于一点 \(t_0 \in J\) 的积分 u 在 J 上可逆，且其逆是方程 \(d\mathbf{x}/dt = -\mathbf{x}\mathbf{a}(t)\) 的解（若 v 是后一方程的取值 e 于点 \(t_0\) 的积分，考虑 uv 与 vu 所满足的线性方程）。由此推出：对每个 \(x_0 \in G\) ，\(d\mathbf{x}/dt = \mathbf{a}(t)\mathbf{x}\) 的取值 \(x_0\) 于点 \(t_0\) 的积分等于 \(\mathbf{u}(t)\mathbf{x}_0\) 。
\item
  设 \(\mathbf{a}(t)\) 、\(\mathbf{b}(t)\) 是 J 上的两个规则函数，u 与 v 分别是方程 \(d\mathbf{x}/dt = \mathbf{a}(t)\mathbf{x}\) 与 \(d\mathbf{x}/dt = \mathbf{x}\mathbf{b}(t)\) 的取值 e 于点 \(t_{0}\) 的积分。证明方程 \(d\mathbf{x}/dt = \mathbf{a}(t)\mathbf{x} + \mathbf{x}\mathbf{b}(t)\) 的取值 \(x_{0}\) 于点 \(t_{0}\) 的积分等于 \(\mathbf{u}(t)\mathbf{x}_{0}\mathbf{v}(t)\) 。
\item
  设 \(\mathbf{a}(t)\) 、\(\mathbf{b}(t)\) 、\(\mathbf{c}(t)\) 、\(\mathbf{d}(t)\) 是 \(\mathbf{J}\) 上的四个规则函数，\((\mathbf{u}, \mathbf{v})\) 是由两个线性方程组成的方程组
\end{enumerate}

\[
\frac {d \mathbf {x}}{d t} = \mathbf {a} (t) \mathbf {x} + \mathbf {b} (t) \mathbf {y}, \quad \frac {d \mathbf {y}}{d t} = \mathbf {c} (t) \mathbf {x} + \mathbf {d} (t) \mathbf {y}.
\]

证明：若 \(\mathbf{v}\) 在 \(J\) 上可逆，则 \(\mathbf{w} = \mathbf{u}\mathbf{v}^{-1}\) 是方程 \(d\mathbf{z}/dt = \mathbf{b}(t) + \mathbf{a}(t)\mathbf{z} - \mathbf{z}\mathbf{d}(t) - \mathbf{z}\mathbf{c}(t)\mathbf{z}\) （“里卡蒂方程”）的一个积分；反之亦然。由此推出：后一方程在 \(J\) 上的每个取值 \(\mathbf{x}_0\) 于点 \(t_0\) 的积分都具有形式 \((A(t)\mathbf{x}_0 + B(t)\mathbf{e})(C(t)\mathbf{x}_0 + D(t)\mathbf{e})^{-1}\) ，其中 \(A(t), B(t), C(t)\) 与 \(D(t)\) 是由 \(J\) 到 \(\mathcal{L}(E)\) 中的连续映射，满足恒等式 \(U.(xy) = (U.x)y\) 。

\begin{enumerate}
\def\labelenumi{\alph{enumi})}
\setcounter{enumi}{3}
\tightlist
\item
  设 \(w_{1}\) 是 \(d\mathbf{z}/dt = \mathbf{b}(t) + \mathbf{a}(t)\mathbf{z} - \mathbf{z}\mathbf{d}(t) - \mathbf{z}\mathbf{c}(t)\mathbf{z}\) 在 J 上的一个积分；证明：若 w 是这个方程的另一个积分，使得 \(w - w_{1}\) 在 J 上可逆，则 \(w - w_{1}\) 可以用方程 \(d\mathbf{x}/dt = -(\mathbf{d} + \mathbf{c}\mathbf{w}_{1})\mathbf{x}\) 与 \(d\mathbf{x}/dt = \mathbf{x}(\mathbf{a} - \mathbf{w}_{1}\mathbf{c})\) 的积分表出（考虑 \((\mathbf{w} - \mathbf{w}_{1})^{-1}\) 所满足的方程，并用 b)）。
\end{enumerate}

\begin{enumerate}
\def\labelenumi{\arabic{enumi})}
\setcounter{enumi}{8}
\tightlist
\item
  设 \(y_{k}\) ( \(1 \leqslant k \leqslant n\) ) 是 \(n\) 个定义在区间 \(I \subset \mathbf{R}\) 上、具有连续 \((n - 1)^{th}\) 导数的实函数。
\end{enumerate}

\begin{enumerate}
\def\labelenumi{\alph{enumi})}
\item
  证明：若 \(n\) 个函数 \(y_{k}\) 线性相关，则矩阵 \((y_{k}^{(h)}(t))_{0\leqslant h\leqslant n - 1, 1\leqslant k\leqslant n}\) 的秩 \(< n\) （在每一点 \(t\in I\) 处）。
\item
  反之，若对每个 \(t \in I\) 矩阵 \((y_{k}^{(h)}(t))_{0 \leqslant h \leqslant n-1, 1 \leqslant k \leqslant n}\) 的秩 \textless{} n，证明：在每个非空开区间 \(J \subset I\) 中存在一个非空开区间 \(U \subset J\) ，使得 \(y_{k}\) 在 U 上的限制线性相关（若 p 是那些数 q \textless{} n ——使函数 \(y_{k}\) 中任意 q 个的朗斯基行列式在 J 上恒为零的数——中最小的一个，考虑一点 \(a \in J\) ，在该处函数 \(y_{k}\) 中 p - 1 个的朗斯基行列式不为零，并证明函数 \(y_{k}\) 中 p 个在 a 的一个邻域上是一个 p - 1 阶线性方程的积分）。
\item
  令 \(y_{1}(t) = t^{2}\) （对 \(t \in \mathbf{R}\) ），\(y_{2}(t) = t^{2}\) （对 \(t \geqslant 0\) ），\(y_{2}(t) = -t^{2}\) （对 \(t < 0\) ）；证明 \(y_{1}\) 与 \(y_{2}\) 在 \(\mathbf{R}\) 上有连续导数，且 \(y_{1}y_{2}' - y_{2}y_{1}' = 0\) ，但 \(y_{1}\) 与 \(y_{2}\) 在 \(\mathbf{R}\) 上不线性相关。
\end{enumerate}

* 10) 设 \(t \mapsto X(t)\) 是一个由区间 \(\mathbf{J} \subset \mathbf{R}\) 到 \(n\) 阶复矩阵空间中的映射。假设导数 \(t \mapsto X'(t)\) 存在且在 \(\mathbf{J}\) 上连续，并且使得 \(X(t)X'(t) = X'(t)X(t)\) 对每个 \(t \in \mathbf{J}\) 成立。

\begin{enumerate}
\def\labelenumi{\alph{enumi})}
\item
  进一步假设对每个 \(t \in J\) ，\(X(t)\) 的特征值互异。证明这时存在一个常数可逆矩阵 \(P_{0}\) ，使得 \(P_{0}X(t)P_{0}^{-1} = D(t)\) ，其中 \(D(t)\) 是一个对角矩阵，从而 \(X(t_{1})\) 与 \(X(t_{2})\) 对 t 在 J 中的每一对值 \(t_{1}, t_{2}\) 可交换。（在 J 的每一点的一个邻域上写出 \(X(t) = P(t)D(t)P(t)^{-1}\) ，并建立 \(P(t)\) 所满足的微分方程。）
\item
  矩阵
\end{enumerate}

\[
X (t) = \left( \begin{array}{c c c} t ^ {2} & t ^ {3} & t ^ {4} \\ - 2 t & - 2 t ^ {2} & - 2 t ^ {3} \\ 1 & t & t ^ {2} \end{array} \right)
\]

与其导数可交换，但 a) 的结论不成立。*

\section*{历史注记}
\phantomsection
\addcontentsline{toc}{section}{历史注记}

（注意：罗马数字指标是指本注末尾的参考文献。）

正如我们已经看到的（第 I、II、III 章的历史注记），导致微分方程积分法的问题属于 XVII \(^{th}\) 世纪无穷小演算的奠基者们（特别是笛卡儿与巴罗）最早处理的问题之列。从那时起，微分方程理论就没有停止过考验数学家们的敏锐，并且一直是应用分析学中各种各样方法的一片沃土；它所引发的问题远未全部解决，因而人们对它的兴趣经久不衰，以至于它构成了数学与实验科学之间最持久、最富有成果的接触点之一：后者常常从中获得宝贵的帮助，而作为交换又不断提供新的问题。

在一部关于微分方程的现代完整研究所应包含的众多章节中，我们在这里只想阐述其中最初步的两章，它们分别处理存在性定理与线性方程，并假定变量为实变量。这也是我们把这段简短的历史叙述限于这两个方面的原因。从 XVIII \(^{th}\) 世纪初起，数学家们就确信，n 阶微分方程的“一般”积分依赖于 n 个任意常数，而且“一般地”存在一个且仅一个积分，它在变量的一个给定值 \(x_{0}\) 处使函数及其前 n - 1 阶导数取给定的值：支撑这一信念的是（追溯到牛顿的）借助微分方程本身及其前 n 个系数逐个计算解在点 \(x_{0}\) 处的泰勒展开式系数的过程。但直到柯西，才有人研究这样得到的级数的收敛性，并证明其和是微分方程的一个解；当然，那时更谈不上非解析微分方程。在柯西为证明微分方程积分的存在性而发明的各种方法中，我们所遵循的那种方法（(IV) 与 (IV bis)），稍后由利普希茨加以推广，对于非解析方程与积分逼近的情形特别有价值。

线性微分方程属于最早引起注意的问题之列。莱布尼茨与雅各布·伯努利用两次求积积分了一阶线性方程（(I), v. II, p.~731）；任意右端的常系数任意线性方程的一般积分由欧拉得到（(II), p.~296-354）；达朗贝尔用同样的方式求解了常系数线性方程组。稍后，拉格朗日 (III) 处理了 \(n\) 阶线性方程的一般理论，认识到齐次方程的积分是 \(n\) 个积分常数的线性函数，引入了伴随方程，发现了当已知若干特解时齐次方程的降阶法，以及求解非齐次方程的常数变易法。拉格朗日理论中的模糊之处（特别是关于积分的线性无关性）由柯西澄清，他的阐述 (IV bis) 除矩阵记号与初等因子理论所带来的细节改进外，几乎已成定论。

\section*{参考文献}
\phantomsection
\addcontentsline{toc}{section}{参考文献}

\begin{enumerate}
\def\labelenumi{(\Roman{enumi})}
\item
  JAKOB BERNOULLI, Opera, 2 vol., Genève (Cramer-Philibert), 1744.
\item
  L. EULER, Opera Omnia: Institutiones calculi integralis, (1) t. XII, Leipzig-Berlin (Teubner), 1914.
\item
  J.-L. LAGRANGE, Œuvres, Paris (Gauthier-Villars), 1867-1890: a) Solution de différents problèmes de calcul intégral, t. I, p.~471; b) Sur le mouvement des nœuds des orbites planétaires, t. IV, p.~111.
\item
  A.-L. CAUCHY, Œuvres complètes, (2), t. XI, Paris (Gauthier-Villars), 1913, p.~399 (= Exercices d'Analyse, Paris, 1840, t. I, p.~327).
\end{enumerate}

(IV bis) A.-L. CAUCHY, dans Leçons de calcul différentiel et de calcul intégral, rédigées principalement d'après les méthodes de M. A.-L. Cauchy, par l'abbé Moigno, t. II, Paris, 1844.

\chapter{函数的局部研究}

\section{滤过集上的函数比较}

设 E 是一个被具有基 F 的滤子所滤过的集合 (Gen.~Top., I, p.~58)；在本章中我们将考虑这样的函数：它们的定义域是属于滤子基 F 的 E 的一个子集（该子集依赖于函数），并且它们的值或者取在实数域 R 中，或者更一般地取在一个赋值域上的赋范向量空间中 (Gen.~Top., IX, p.~170)。

在应用中，E 大多是实空间 \(R^{n}\) 的一个子空间，或者是扩充实直线 \(\overline{R}\) ，而 \(F\) 则是 E 的某个闭包点的邻域滤子在 E 上的迹，或者是 E 的相对紧子集的补集所成的滤子（“无穷远点的邻域”）。

一般说来，仅知道这样一个函数沿 \(\mathfrak{F}\) 趋于一个给定的极限，并不足以处理由这个函数构成的表达式的所有“沿 \(\mathfrak{F}\) 过渡到极限”的问题。

例如，当实变量 \(x\) 趋于 \(+\infty\) 时，三个函数 \(x\) 、\(x^2\) 与 \(\sqrt{x}\) 都趋于 \(+\infty\) ，但在诸表达式

\[
(x + 1) ^ {2} - x ^ {2}, \qquad (x + 1) - x, \quad \sqrt {x + 1} - \sqrt {x},
\]

中，第一个趋于 \(+\infty\) ，第二个趋于 1，第三个趋于 0。

因此，重要的不只是知道函数沿 \(\mathfrak{F}\) 的极限值（当这极限存在时），还要知道函数接近其极限的“方式”；换句话说，我们要对趋于同一极限的函数所成之集进行分类。
